\documentclass[a4paper, 12pt]{article}
\usepackage{tcolorbox}
\usepackage{a4wide,amssymb,epsfig,latexsym,multicol,array,hhline,fancyhdr}
\usepackage{longtable}
\usepackage{booktabs}
\usepackage{amsmath}
\usepackage{lastpage}
\usepackage{placeins}
\usepackage{subfig}
\usepackage[lined,boxed,commentsnumbered]{algorithm2e}
\usepackage{enumerate}
\usepackage{color}
\usepackage{graphicx}			
\usepackage{ragged2e}
\usepackage{array}
\usepackage{tabularx, caption}
\usepackage{multirow}
\usepackage[framemethod=tikz]{mdframed}
\usepackage{multicol}
\usepackage{rotating}
\usepackage{graphics}
\usepackage{geometry}
\usepackage{setspace}
\usepackage{epsfig}
\usepackage{tikz}
\usepackage{makecell}
\usepackage[utf8]{vietnam}
\usepackage{listings}
\usepackage[table]{xcolor}
\usetikzlibrary{arrows, snakes, backgrounds, fit, calc, positioning, arrows.meta, shapes.geometric}
\usepackage{hyperref}
\usepackage{indentfirst}
\usepackage{float}
\usepackage{xltabular}
\def\thesislayout{	% A4: 210 × 297
	\geometry{
		paper=a4paper,
		total={160mm,240mm},  % fix over page
		left=2cm,
        right=2cm,
		top=1.5cm,
        bottom=1.5cm,
        includeheadfoot=true,
        headheight=30pt
	}
}
\hypersetup{urlcolor=blue,linkcolor=black,citecolor=black,colorlinks=true}
\thesislayout
\AtBeginDocument{\renewcommand*\contentsname{Mục lục}}
\setlength{\headheight}{40pt}
\pagestyle{fancy}
\fancyhead{}
\fancyhead[L]{
 \begin{tabular}{rl}
    \begin{picture}(25,15)(0,0)
    \put(-3,-9){\includegraphics[width=17mm]{logobk.png}}
    %\put(0,-8){\epsfig{width=10mm,figure=hcmut.eps}}
   \end{picture}&
	%\includegraphics[width=8mm, height=8mm]{hcmut.png} & %
	\begin{tabular}{l}
		\textbf{\bf \ttfamily Trường Đại học Bách khoa - ĐHQG-HCM}\\
		\textbf{\bf \ttfamily Khoa Khoa học và Kỹ thuật Máy tính}
	\end{tabular} 	
 \end{tabular}
}
\fancyhead[R]{
	\begin{tabular}{l}
		\tiny \bf \\
		\tiny \bf 
	\end{tabular}  }
\fancyfoot{} % clear all footer fields
\fancyfoot[L]{\scriptsize \texttt{Bài tập lớn - Kho dữ liệu và Hệ hỗ trợ quyết định (CO4031) - Năm học 2025-2026}}
\fancyfoot[R]{\scriptsize Trang {\thepage}/\pageref{LastPage}}
\renewcommand{\headrulewidth}{0.3pt}
\renewcommand{\footrulewidth}{0.3pt}
\renewcommand{\contentsname}{Mục lục}
%%%
\setcounter{secnumdepth}{4}
\setcounter{tocdepth}{3}
\makeatletter
\newcounter {subsubsubsection}[subsubsection]
\renewcommand\thesubsubsubsection{\thesubsubsection .\@alph\c@subsubsubsection}
\newcommand\subsubsubsection{\@startsection{subsubsubsection}{4}{\z@}%
                                     {-3.25ex\@plus -1ex \@minus -.2ex}%
                                     {1.5ex \@plus .2ex}%
                                     {\normalfont\normalsize\bfseries}}
\newcommand*\l@subsubsubsection{\@dottedtocline{3}{10.0em}{4.1em}}
\newcommand*{\subsubsubsectionmark}[1]{}
\makeatother

\definecolor{dkgreen}{rgb}{0,0.6,0}
\definecolor{gray}{rgb}{0.5,0.5,0.5}
\definecolor{mauve}{rgb}{0.58,0,0.82}
\definecolor{codeblue}{rgb}{0.0,0.0,255}     % Keywords
\definecolor{codegreen}{rgb}{0,0.6,0}        % Comments
\definecolor{codepurple}{rgb}{138,0,226}   % Strings
\definecolor{codeorange}{rgb}{0.8,0.33,0}    % Types
\definecolor{codegray}{rgb}{0.255,0.0,0.255}     % Line numbers
\definecolor{codepreprocessor}{rgb}{0.5,0,0} % Preprocessor commands
\definecolor{codebackground}{rgb}{0.95,0.95,0.92} % Background color

\lstset{
    language=R,
    basicstyle=\ttfamily\scriptsize,
    keywordstyle=\bfseries\color{codeblue},
    commentstyle=\color{codegreen},
    stringstyle=\color{codepurple},
    showstringspaces=false,
    numbers=left,
    numberstyle=\tiny\color{gray},
    breaklines=true,
    frame=single,
    rulecolor=\color{black} % Viền khung sẽ hiển thị màu đen
}
\thesislayout
\begin{document}
\onehalfspacing
\renewcommand{\refname}{Tài liệu tham khảo}

\begin{titlepage}
\begin{tikzpicture}[remember picture,overlay,inner sep=0,outer sep=0]
     \draw[blue!70!black,line width=4pt] ([xshift=-1.5cm,yshift=-2cm]current page.north east) coordinate (A)--([xshift=1.5cm,yshift=-2cm]current page.north west) coordinate(B)--([xshift=1.5cm,yshift=2cm]current page.south west) coordinate (C)--([xshift=-1.5cm,yshift=2cm]current page.south east) coordinate(D)--cycle;

     \draw ([yshift=0.5cm,xshift=-0.5cm]A)-- ([yshift=0.5cm,xshift=0.5cm]B)--
     ([yshift=-0.5cm,xshift=0.5cm]B) --([yshift=-0.5cm,xshift=-0.5cm]B)--([yshift=0.5cm,xshift=-0.5cm]C)--([yshift=0.5cm,xshift=0.5cm]C)--([yshift=-0.5cm,xshift=0.5cm]C)-- ([yshift=-0.5cm,xshift=-0.5cm]D)--([yshift=0.5cm,xshift=-0.5cm]D)--([yshift=0.5cm,xshift=0.5cm]D)--([yshift=-0.5cm,xshift=0.5cm]A)--([yshift=-0.5cm,xshift=-0.5cm]A)--([yshift=0.5cm,xshift=-0.5cm]A);


     \draw ([yshift=-0.3cm,xshift=0.3cm]A)-- ([yshift=-0.3cm,xshift=-0.3cm]B)--
     ([yshift=0.3cm,xshift=-0.3cm]B) --([yshift=0.3cm,xshift=0.3cm]B)--([yshift=-0.3cm,xshift=0.3cm]C)--([yshift=-0.3cm,xshift=-0.3cm]C)--([yshift=0.3cm,xshift=-0.3cm]C)-- ([yshift=0.3cm,xshift=0.3cm]D)--([yshift=-0.3cm,xshift=0.3cm]D)--([yshift=-0.3cm,xshift=-0.3cm]D)--([yshift=0.3cm,xshift=-0.3cm]A)--([yshift=0.3cm,xshift=0.3cm]A)--([yshift=-0.3cm,xshift=0.3cm]A);

   \end{tikzpicture}
\begin{center}
\textbf{ĐẠI HỌC QUỐC GIA THÀNH PHỐ HỒ CHÍ MINH}\\
\textbf{TRƯỜNG ĐẠI HỌC BÁCH KHOA}\\
\textbf{KHOA KHOA HỌC VÀ KỸ THUẬT MÁY TÍNH} \\
\end{center}
\begin{figure}[h!]
\begin{center}
\includegraphics[width=8cm]{logobk.png}
\end{center}
\end{figure}
\begin{center}
\begin{tabular}{c}
	\multicolumn{1}{c}{\textbf{{\Large BÁO CÁO BÀI TẬP LỚN}}}\\
    \multicolumn{1}{c}{\textbf{{\Large MÔN: Kho dữ liệu và hệ hỗ trợ quyết định (CO4031)}}}\\
	~~\\
	\hline
	\\
	\multicolumn{1}{c}{\textbf{\Large \parbox{\textwidth}{\centering BÀI TẬP LỚN Kho dữ liệu và hệ hỗ trợ quyết định}}}\\
	\\
	\hline
\end{tabular}
\end{center}
\begin{table}[h]
    \centering
    \begin{tabular}{rrrrlr}
        & \textbf{Giảng viên hướng dẫn:} && & ThS. Lê Thị Bảo Thu & \\
        & \textbf{Lớp:} &&& L01 & \\
        & \textbf{Nhóm:} &&& 4 & \\
        
        & \textbf{Sinh viên thực hiện:} &&& Nguyễn Trọng Nhân  & 2212370 \\
        &                      &&& Nguyễn Phú Tuệ     & 2213813 \\        
        &                      &&& Nguyễn Hoàng Minh       & 2212057 \\
        &                      &&& Hoàng Giữ Tiến Nhất & 2312460 \\
        &                      &&& Mai Xuân Phúc        & 2312687 \\
    \end{tabular}
\end{table}
% \vspace{0.3cm}
\begin{center}
\textit{Thành phố Hồ Chí Minh, Tháng 4/2026}
\end{center}
\end{titlepage}

%%%%%%%%%%%%%%%%%%%%%%%%%%%%%%%%%%%%%%%%%%%%%%
\newpage
\begin{center}
    \Large{\textbf{Danh sách thành viên \& khối lượng công việc}}
\end{center}
\begin{table}[h]
    \centering
    \renewcommand{\arraystretch}{1.3} % Điều chỉnh khoảng cách giữa các dòng
    \resizebox{\textwidth}{!}{%
        \begin{tabular}{|c|c|l|l|l|c|}
            \hline
            \textbf{STT} & \textbf{MSSV} & \textbf{Họ} & \textbf{Tên} & \textbf{Nhiệm vụ} & \textbf{\% Hoàn thành} \\
            \hline
            1 & 2212370  & Nguyễn Trọng    & Nhân & Mục 1, Mục 2.1, 2.2, Kiểm thử dữ liệu & 100\% \\
            \hline
            2 & 2213813  & Nguyễn Phú  & Tuệ  & Mục 2.3, 2.4, Thiết kế Star Schema \& OLAP & 100\% \\
            \hline
            3 & 2212057  & Nguyễn Hoàng  & Minh  & Mục 3 (ETL Bronze -- Silver -- Gold), Dashboard & 100\% \\
            \hline
            4 & 2312460  & Hoàng Giữ Tiến  & Nhất & Kiến trúc DW, Mục 3.3.4 (Incremental), Mục 4 (DSS) & 100\% \\
            \hline
            5 & 2312687  & Mai Xuân  & Phúc  & Mục 4 (Thực nghiệm DSS), Mục 5, Mục 6 & 100\% \\
            \hline
        \end{tabular}%
    }
\end{table}
\newpage

\section*{Lời cảm ơn}
Nhóm thực hiện xin gửi lời cảm ơn chân thành và sâu sắc đến giảng viên hướng dẫn môn học Kho dữ liệu và Hệ hỗ trợ quyết định (CO4031), Khoa Khoa học và Kỹ thuật Máy tính, Trường Đại học Bách Khoa -- ĐHQG TP.HCM. Những kiến thức nền tảng và định hướng phương pháp luận về mô hình hóa dữ liệu đa chiều, thiết kế quy trình trích xuất -- biến đổi -- nạp (ETL), khai thác truy vấn phân tích trực tuyến (OLAP) và xây dựng hệ hỗ trợ quyết định (DSS) được giảng dạy trong học kỳ là cơ sở quan trọng giúp nhóm hoàn thành đề tài nghiên cứu và xây dựng hệ thống thực nghiệm này.

\newpage
\section*{Tóm tắt}
Báo cáo trình bày quá trình phân tích, thiết kế và hiện thực hóa hệ thống Kho dữ liệu (Data Warehouse) tích hợp Hệ hỗ trợ quyết định (Decision Support System -- DSS) cho bài toán bán lẻ đa quốc gia trên bộ dữ liệu Global Superstore gồm 51.290 dòng giao dịch tại 147 quốc gia trong giai đoạn 2011--2014. Hệ thống được triển khai theo kiến trúc phân tầng Medallion (\texttt{bronze}, \texttt{silver}, \texttt{gold}) trên hệ quản trị cơ sở dữ liệu PostgreSQL 17 đóng gói trong Docker. Quy trình ETL tự động hóa hoàn toàn việc kiểm toán tệp nguồn qua mã băm SHA-256, xử lý ngày tháng đa định dạng, điền khuyết mã bưu chính, đảm bảo tính lũy đẳng và hỗ trợ nạp gia tăng (Incremental Loading) thông qua cơ chế bắt thay đổi dữ liệu. Tại tầng Gold, mô hình hình sao (Star Schema) với bảng sự kiện \texttt{gold.fact\_sales} và 5 bảng chiều chuẩn hóa khóa thay thế đạt độ chính xác đối soát tài chính tuyệt đối (12.642.501,91 USD doanh thu, 1.467.457,29 USD lợi nhuận), kết hợp cùng 7 kịch bản truy vấn OLAP và 3 bảng Materialized Views cho thời gian phản hồi dưới 50 ms. Trên nền tảng kho dữ liệu, phân hệ DSS triển khai ba mô-đun gồm: dự báo chuỗi thời gian bằng thuật toán San bằng mũ có xu thế (MAPE đạt 14,72\%), khai phá luật kết hợp giỏ hàng Apriori hai cấp độ trên 25.035 đơn hàng (phát hiện 14 luật cấp tiểu ngành hàng và 27 luật cấp sản phẩm với Lift lên tới 315,57), và mô phỏng kịch bản What-If dựa trên độ co giãn của cầu theo giá. Kết quả toàn hệ thống được tích hợp trực tiếp lên dịch vụ FastAPI, giao diện Web tương tác Streamlit và báo cáo quản trị Power BI Desktop.

\newpage
\tableofcontents
\newpage
\renewcommand{\listfigurename}{Danh mục hình ảnh}
\listoffigures
\addcontentsline{toc}{section}{\listfigurename}
\newpage
%%%%%%%%%%%%%%%%%%%%%%%%%%%%%%%%%%%%%%%%%%%%%
\section{Giới thiệu đề tài}

\subsection{Bối cảnh và vấn đề}

Trong bối cảnh thương mại điện tử và bán lẻ đa quốc gia phát triển mạnh mẽ, các doanh nghiệp phân phối hàng hóa toàn cầu phải vận hành đồng thời trên hàng trăm thị trường quốc gia với danh mục sản phẩm lên đến hàng chục nghìn mặt hàng (Stock Keeping Units -- SKU). Hoạt động giao dịch hàng ngày tạo ra một khối lượng lớn dữ liệu vận hành phát sinh từ các hệ thống xử lý giao dịch trực tuyến (Online Transaction Processing -- OLTP), bao gồm thông tin đơn đặt hàng, lịch trình giao vận, phân khúc khách hàng, cấu trúc giá bán, chính sách chiết khấu và lợi nhuận ròng trên từng dòng sản phẩm. Tuy nhiên, việc lưu trữ dữ liệu dưới dạng các tệp tin phẳng rời rạc hoặc trong các cơ sở dữ liệu giao dịch chưa qua chuẩn hóa khiến doanh nghiệp gặp nhiều rào cản khi cần tổng hợp báo cáo phân tích lịch sử dài hạn và hỗ trợ ra quyết định ở cấp độ chiến lược.

Qua khảo sát bài toán thực tế tại doanh nghiệp bán lẻ toàn cầu Global Superstore với hơn 51.290 bản ghi giao dịch chi tiết trải rộng trên 147 quốc gia trong giai đoạn 2011--2014, nhóm thực hiện nhận diện ba vấn đề cốt lõi cần giải quyết. Thứ nhất, dữ liệu nguồn xuất phát từ các hệ thống vận hành có độ phân mảnh và không đồng nhất cao về mặt định dạng, điển hình là sự pha trộn giữa nhiều quy chuẩn ngày tháng, tình trạng khuyết thiếu mã bưu chính (Postal Code) tại các đơn hàng ngoài lãnh thổ Hoa Kỳ và sự trùng lặp thông tin danh mục khách hàng cũng như địa lý qua hàng chục nghìn dòng giao dịch. Nếu truy vấn trực tiếp trên nguồn dữ liệu thô này, chi phí tính toán sẽ rất lớn và dễ dẫn đến sai lệch số liệu tài chính khi tổng hợp theo nhiều chiều.

Thứ hai, trong công tác quản trị hiệu quả kinh doanh, ban điều hành đối mặt với bài toán suy giảm biên lợi nhuận cục bộ do chính sách chiết khấu thiếu kiểm soát. Mặc dù tổng doanh thu toàn cầu tăng trưởng đều đặn qua từng năm và đạt ngưỡng 12,64 triệu USD, nhiều khu vực thị trường (như Trung Đông -- EMEA hoặc một số quốc gia tại Châu Phi và Đông Nam Á) và một số phân nhóm ngành hàng (điển hình là nhóm Bàn làm việc -- Tables thuộc ngành Nội thất) lại ghi nhận mức lợi nhuận ròng âm nặng nề khi áp dụng mức chiết khấu từ 30\% đến trên 50\%. Các báo cáo thống kê tĩnh truyền thống chỉ dừng lại ở việc mô tả con số tổng hợp mà không cung cấp khả năng khoan sâu (Drill-down) hoặc xoay chiều (Pivot) linh hoạt để bóc tách nguyên nhân gây thất thoát dòng tiền.

Thứ ba, các nhà quản lý thiếu một hệ thống hỗ trợ quyết định (Decision Support System -- DSS) mang tính định lượng để dự báo nhu cầu thị trường và đánh giá rủi ro trước khi ban hành chính sách. Khi đứng trước các câu hỏi quản trị như cần chuẩn bị ngân sách và lượng hàng tồn kho bao nhiêu cho hai quý tiếp theo, nên ghép cặp những sản phẩm nào vào cùng một gói khuyến mãi bán chéo (Cross-selling) để gia tăng giá trị giỏ hàng trung bình, hay việc điều chỉnh tăng/giảm 5\% tỷ lệ chiết khấu sẽ tác động như thế nào đến tổng lợi nhuận dựa trên độ co giãn của cầu theo giá, việc ra quyết định dựa trên cảm tính kinh nghiệm thường tiềm ẩn rủi ro tài chính cao. Những vấn đề thực tiễn này đặt ra yêu cầu cấp thiết về việc xây dựng một kho dữ liệu (Data Warehouse) tập trung, chuẩn hóa theo mô hình đa chiều kết hợp cùng các mô hình phân tích dự báo và khai phá dữ liệu tự động.

\subsection{Mục tiêu thực hiện}

Đề tài hướng tới mục tiêu thiết kế, xây dựng và kiểm chứng một hệ thống Kho dữ liệu (Data Warehouse) tích hợp Hệ hỗ trợ quyết định (Decision Support System) hoàn chỉnh từ khâu tiếp nhận dữ liệu thô đến khâu trực quan hóa và khuyến nghị hành động cho bài toán Global Superstore. Ở tầng nền tảng dữ liệu, hệ thống cần thiết lập một quy trình trích xuất, biến đổi và nạp dữ liệu (Extract -- Transform -- Load, viết tắt là ETL) tự động theo kiến trúc phân tầng Medallion (Bronze -- Silver -- Gold) trên hệ quản trị cơ sở dữ liệu quan hệ PostgreSQL. Quy trình này phải đảm bảo tính lũy đẳng (Idempotency), khả năng lưu vết kiểm toán từng đợt nạp thông qua mã băm nội dung SHA-256, khả năng nạp gia tăng (Incremental Loading) đối với các lô dữ liệu giao dịch mới phát sinh mà không cần xử lý lại dữ liệu lịch sử, đồng thời bảo toàn tuyệt đối độ chính xác của các chỉ số tài chính giữa các tầng dữ liệu.

Ở tầng lưu trữ và truy vấn phân tích trực tuyến (Online Analytical Processing -- OLAP), mục tiêu của đề tài là chuyển đổi cấu trúc dữ liệu bảng phẳng ban đầu thành mô hình hình sao (Star Schema) tuân thủ phương pháp luận của Ralph Kimball. Mô hình cần tách biệt rõ ràng bảng sự kiện trung tâm (\texttt{gold.fact\_sales}) lưu trữ các độ đo định lượng với năm bảng chiều (\texttt{dim\_date}, \texttt{dim\_customer}, \texttt{dim\_product}, \texttt{dim\_geography}, \texttt{dim\_ship\_mode}) sử dụng khóa thay thế số nguyên (Surrogate Keys). Trên nền tảng mô hình hình sao này, hệ thống phải hiện thực hóa đầy đủ các phép toán OLAP cốt lõi (Roll-up, Drill-down, Slice, Dice, Pivot, CUBE và các hàm cửa sổ phân tích) cùng hệ thống các bảng tổng hợp vật lý (Materialized Views) nhằm tối ưu hóa thời gian phản hồi truy vấn phục vụ bảng điều khiển quản trị.

Ở tầng hỗ trợ ra quyết định (DSS), mục tiêu trọng tâm là phát triển ba mô-đun phân tích chuyên sâu có khả năng vận hành trực tiếp trên dữ liệu sạch từ tầng Gold và lưu trữ kết quả đầu ra vào lược đồ chuyên biệt (\texttt{dss}). Mô-đun thứ nhất thực hiện phân tích dự báo chuỗi thời gian (Predictive Analytics) để ước lượng doanh số bán hàng và sản lượng tiêu thụ trong tương lai kèm theo khoảng tin cậy thống kê 95\%, đồng thời đánh giá khách quan chất lượng dự báo qua các chỉ số sai số thực nghiệm MAE, RMSE và MAPE. Mô-đun thứ hai ứng dụng khai phá luật kết hợp (Association Rule Mining) trên hàng chục nghìn giỏ hàng thực tế để phát hiện các mối tương quan mua sắm đồng thời giữa các mặt hàng và tiểu ngành hàng dựa trên các độ đo Support, Confidence và Lift. Mô-đun thứ ba xây dựng bộ máy mô phỏng kịch bản What-If dựa trên mô hình kinh tế lượng về độ co giãn của cầu theo giá, kết hợp cùng giao diện ứng dụng Web tương tác (Streamlit, FastAPI) và báo cáo Power BI để cung cấp các khuyến nghị điều hành trực quan cho người ra quyết định.

\subsection{Phạm vi và phương pháp tiếp cận}

Về phạm vi dữ liệu và nghiệp vụ, đề tài tập trung khai thác bộ dữ liệu giao dịch lịch sử của Global Superstore gồm 51.290 dòng chi tiết đơn hàng (tương ứng với 25.035 đơn đặt hàng độc lập của 1.590 khách hàng đối với 10.292 mã sản phẩm) trong giai đoạn 48 tháng liên tục từ tháng 01/2011 đến hết tháng 12/2014, cùng một tập dữ liệu kiểm thử nạp gia tăng gồm 500 dòng giao dịch phát sinh trong Quý 1/2015. Phạm vi phân tích bao quát toàn bộ bảy khu vực thị trường quốc tế (APAC, EU, US, LATAM, EMEA, Africa, Canada) và ba ngành hàng chính (Technology, Furniture, Office Supplies) với 17 phân nhóm sản phẩm con. Các chiều phân tích tập trung vào hiệu quả tài chính (doanh thu, lợi nhuận ròng, biên lợi nhuận), chính sách chiết khấu, chi phí và thời gian vận chuyển, hành vi mua sắm theo phân khúc khách hàng và cấu trúc giỏ hàng.

Về phương pháp tiếp cận kỹ thuật, nhóm thực hiện áp dụng quy trình phát triển hệ thống dữ liệu theo hướng mô-đun hóa và hướng cấu hình. Trong giai đoạn phân tích và thiết kế kho dữ liệu, phương pháp thiết kế bốn bước của Kimball được áp dụng chặt chẽ: xác định quy trình nghiệp vụ bán hàng làm trọng tâm, tuyên bố độ mịn (Grain) của bảng sự kiện ở mức độ chi tiết nhất là mỗi dòng sản phẩm trong một đơn hàng (\texttt{row\_id}), thiết kế các chiều phân tích chuẩn hóa và xác định các độ đo cộng được (Additive Measures). Để giải quyết bài toán làm sạch dữ liệu và đồng bộ luồng xử lý, kiến trúc Medallion ba tầng được triển khai kết hợp giữa ngôn ngữ Python (thư viện Pandas, SQLAlchemy, Psycopg2) ở khâu điều phối và sức mạnh xử lý tập hợp của ngôn ngữ SQL chuẩn PostgreSQL ở khâu biến đổi dữ liệu trong cơ sở dữ liệu (In-database ELT/ETL).

Đối với phần phát triển Hệ hỗ trợ quyết định, nhóm tiếp cận theo chu trình bốn giai đoạn của Herbert Simon gồm Thám sát (Intelligence), Thiết kế (Design), Lựa chọn (Choice) và Thực thi (Implementation). Các thuật toán dự báo san bằng mũ Holt-Winters có xu thế, thuật toán khai phá tập phổ biến Apriori và mô hình kinh tế lượng co giãn giá được cài đặt thành các gói chương trình Python độc lập trong thư mục \texttt{src/dss/}. Toàn bộ thông số cấu hình lần chạy, chỉ số đánh giá mô hình và kết quả đầu ra đều được ghi nhận bền vững xuống cơ sở dữ liệu PostgreSQL trước khi cung cấp qua giao diện lập trình ứng dụng RESTful API (FastAPI), giao diện tương tác Streamlit và bộ báo cáo quản trị Power BI Desktop. Cách tiếp cận này đảm bảo tính tách biệt rõ ràng giữa tầng lưu trữ dữ liệu, tầng xử lý thuật toán và tầng trình diễn, giúp hệ thống dễ dàng kiểm thử, bảo trì và mở rộng.

\newpage
\section{Dataset và phân tích nghiệp vụ}

\subsection{Giới thiệu nguồn dữ liệu và cấu trúc dữ liệu}

Nguồn dữ liệu chính được sử dụng trong đề tài là bộ dữ liệu bán lẻ toàn cầu Global Superstore (\texttt{dataset/superstore\_dataset2011-2015.csv}), phản ánh toàn bộ hoạt động đặt hàng và giao nhận của một tập đoàn thương mại điện tử đa quốc gia trong suốt bốn năm tài chính từ ngày 01/01/2011 đến ngày 31/12/2014. Tập dữ liệu gốc bao gồm chính xác 51.290 dòng bản ghi (rows) và 24 cột thuộc tính (columns), với dung lượng lưu trữ vật lý khoảng 12,03 MB và được mã hóa theo bảng mã ký tự Latin-1 (ISO-8859-1 / Windows-1252) do chứa nhiều tên riêng khách hàng và địa danh bằng tiếng Pháp, tiếng Tây Ban Nha và tiếng Đức. Bên cạnh tệp dữ liệu lịch sử chính, hệ thống còn sử dụng tệp dữ liệu gia tăng \texttt{dataset/incremental\_orders\_2015\_q1.csv} gồm 500 bản ghi giao dịch phát sinh trong Quý 1/2015 (có mã định danh \texttt{Row ID} từ 60001 đến 60500) để phục vụ kiểm thử quy trình nạp bổ sung.

Về mặt ngữ nghĩa nghiệp vụ, mỗi dòng bản ghi trong tập dữ liệu đại diện cho một mục hàng chi tiết thuộc một đơn đặt hàng cụ thể (Order Line Item). Một đơn đặt hàng (\texttt{Order ID}) có thể bao gồm một hoặc nhiều dòng bản ghi tương ứng với các mặt hàng khác nhau được khách hàng thanh toán trong cùng một giỏ hàng. Trong quy trình kiểm tra sơ bộ tại tầng tiếp nhận (\texttt{src/etl/bronze/extract.py}), 24 thuộc tính của tập dữ liệu được phân chia thành nhóm 20 thuộc tính bắt buộc (\texttt{REQUIRED\_COLUMNS}) và 4 thuộc tính tùy chọn hoặc bổ trợ (\texttt{OPTIONAL\_COLUMNS}) gồm \texttt{Postal Code}, \texttt{Market}, \texttt{Shipping Cost} và \texttt{Order Priority}. Cấu trúc chi tiết của 24 trường dữ liệu trong nguồn gốc được tổng hợp trong Bảng \ref{tab:raw_schema}.

\begin{table}[H]
\centering
\caption{Cấu trúc các trường thuộc tính trong bộ dữ liệu nguồn Global Superstore}
\label{tab:raw_schema}
\footnotesize
\renewcommand{\arraystretch}{1.22}
\begin{tabularx}{\textwidth}{|l|l|l|X|}
\hline
\textbf{Nhóm nghiệp vụ} & \textbf{Tên thuộc tính gốc} & \textbf{Kiểu dữ liệu đích} & \textbf{Ý nghĩa và đặc điểm dữ liệu} \\
\hline
\multirow{6}{*}{Đơn hàng \& Vận chuyển} 
 & \texttt{Row ID} & \texttt{INTEGER} & Mã định danh duy nhất cho từng dòng chi tiết đơn hàng (1 đến 51290). \\ \cline{2-4}
 & \texttt{Order ID} & \texttt{VARCHAR(50)} & Mã đơn đặt hàng giao dịch (có 25.035 mã đơn hàng phân biệt). \\ \cline{2-4}
 & \texttt{Order Date} & \texttt{DATE} & Ngày khách hàng xác nhận đặt hàng (từ 01/01/2011 đến 31/12/2014). \\ \cline{2-4}
 & \texttt{Ship Date} & \texttt{DATE} & Ngày xuất kho giao hàng thực tế (luôn lớn hơn hoặc bằng \texttt{Order Date}). \\ \cline{2-4}
 & \texttt{Ship Mode} & \texttt{VARCHAR(50)} & Phương thức vận chuyển: Standard Class, Second Class, First Class, Same Day. \\ \cline{2-4}
 & \texttt{Order Priority} & \texttt{VARCHAR(20)} & Mức độ ưu tiên xử lý đơn hàng: Low, Medium, High, Critical. \\
\hline
\multirow{3}{*}{Khách hàng} 
 & \texttt{Customer ID} & \texttt{VARCHAR(50)} & Mã định danh tự nhiên của khách hàng (1.590 mã khách hàng duy nhất). \\ \cline{2-4}
 & \texttt{Customer Name} & \texttt{VARCHAR(100)} & Họ và tên đầy đủ của khách hàng đặt mua. \\ \cline{2-4}
 & \texttt{Segment} & \texttt{VARCHAR(50)} & Phân khúc khách hàng: Consumer (Cá nhân), Corporate (Doanh nghiệp), Home Office. \\
\hline
\multirow{6}{*}{Địa lý \& Thị trường} 
 & \texttt{City} & \texttt{VARCHAR(100)} & Tên thành phố nhận hàng của đơn giao dịch. \\ \cline{2-4}
 & \texttt{State} & \texttt{VARCHAR(100)} & Tên bang, tỉnh hoặc vùng hành chính trực thuộc quốc gia. \\ \cline{2-4}
 & \texttt{Country} & \texttt{VARCHAR(100)} & Tên quốc gia nhận hàng (bao phủ 147 quốc gia trên thế giới). \\ \cline{2-4}
 & \texttt{Postal Code} & \texttt{VARCHAR(20)} & Mã bưu chính khu vực (chỉ có giá trị tại Hoa Kỳ, khuyết thiếu ở các nước khác). \\ \cline{2-4}
 & \texttt{Market} & \texttt{VARCHAR(50)} & Khu vực thị trường vĩ mô: APAC, EU, US, LATAM, EMEA, Africa, Canada. \\ \cline{2-4}
 & \texttt{Region} & \texttt{VARCHAR(50)} & Vùng địa lý chi tiết trực thuộc thị trường (13 vùng kinh doanh). \\
\hline
\multirow{4}{*}{Sản phẩm} 
 & \texttt{Product ID} & \texttt{VARCHAR(50)} & Mã định danh sản phẩm (10.292 mã SKU phân biệt). \\ \cline{2-4}
 & \texttt{Category} & \texttt{VARCHAR(50)} & Ngành hàng chính: Technology, Furniture, Office Supplies. \\ \cline{2-4}
 & \texttt{Sub-Category} & \texttt{VARCHAR(50)} & Tiểu ngành hàng chi tiết (17 phân nhóm sản phẩm con). \\ \cline{2-4}
 & \texttt{Product Name} & \texttt{VARCHAR(255)} & Tên thương mại đầy đủ và mô tả đặc tính của sản phẩm. \\
\hline
\multirow{5}{*}{Độ đo tài chính} 
 & \texttt{Sales} & \texttt{NUMERIC(12,4)} & Doanh thu thuần của dòng sản phẩm (đơn vị: USD, giá trị dương). \\ \cline{2-4}
 & \texttt{Quantity} & \texttt{INTEGER} & Số lượng đơn vị sản phẩm được đặt mua trong dòng giao dịch. \\ \cline{2-4}
 & \texttt{Discount} & \texttt{NUMERIC(5,4)} & Tỷ lệ chiết khấu áp dụng trên dòng hàng (nhận giá trị từ 0,00 đến 0,85). \\ \cline{2-4}
 & \texttt{Profit} & \texttt{NUMERIC(12,4)} & Lợi nhuận ròng thu được (USD, có thể mang giá trị dương hoặc âm). \\ \cline{2-4}
 & \texttt{Shipping Cost} & \texttt{NUMERIC(12,4)} & Chi phí vận chuyển phát sinh cho dòng sản phẩm tương ứng (USD). \\
\hline
\end{tabularx}
\end{table}

\subsection{Exploratory Data Analysis (EDA)}

Quá trình phân tích khám phá dữ liệu (Exploratory Data Analysis -- EDA) được thực hiện toàn diện trên 51.290 bản ghi nhằm làm rõ quy mô hoạt động, xu hướng biến động theo thời gian, cơ cấu đóng góp theo không gian địa lý và hiệu quả sinh lời của từng danh mục sản phẩm. Tính trên toàn bộ giai đoạn 2011--2014, doanh nghiệp đạt tổng doanh thu lũy kế là 12.642.501,91 USD, tổng lợi nhuận ròng đạt 1.467.457,29 USD (tương ứng với biên lợi nhuận gộp toàn cục là 11,61\%), tổng sản lượng hàng hóa bán ra đạt 178.312 đơn vị sản phẩm trên 25.035 đơn đặt hàng, với tỷ lệ chiết khấu trung bình là 14,29\%.

Xét theo chiều thời gian qua bốn năm tài chính, hoạt động kinh doanh của Global Superstore thể hiện xu hướng tăng trưởng mở rộng quy mô rõ rệt và bền vững. Năm 2011, hệ thống ghi nhận 4.440 đơn hàng (8.998 dòng giao dịch) với doanh thu 2.259.450,90 USD và lợi nhuận 248.940,81 USD (biên lợi nhuận 11,02\%). Sang năm 2012, số lượng đơn hàng tăng lên 5.343 đơn (10.962 dòng), mang về 2.677.438,69 USD doanh thu và 307.415,28 USD lợi nhuận (biên lợi nhuận 11,48\%). Đà tăng trưởng tiếp tục bứt phá trong năm 2013 với 6.721 đơn hàng (13.799 dòng), đạt doanh thu 3.405.746,45 USD và lợi nhuận 406.935,23 USD (biên lợi nhuận đạt đỉnh 11,95\%). Đến năm 2014, quy mô giao dịch đạt 8.531 đơn hàng (17.531 dòng), đưa doanh thu lên mức 4.299.865,87 USD và lợi nhuận đạt 504.165,97 USD (biên lợi nhuận 11,73\%). Như vậy, chỉ sau bốn năm vận hành, doanh số hàng năm của doanh nghiệp đã tăng gần gấp đôi (+90,31\% từ năm 2011 đến năm 2014). Khi phân tích sâu theo chuỗi 48 tháng, dữ liệu bộc lộ tính chu kỳ mùa vụ (Seasonality) rất rõ nét: doanh số luôn ở mức thấp nhất vào tháng 1 và tháng 2 đầu năm (dao động từ 91.152 USD đến 98.898 USD trong năm 2011), sau đó tăng tốc mạnh vào tháng 6 và đạt đỉnh cao nhất vào quý 4 hàng năm (đặc biệt là tháng 11 và tháng 12/2014 với doanh thu lần lượt đạt 555.279,03 USD và 503.143,69 USD).

Xét theo chiều không gian địa lý và thị trường tiêu thụ (Bảng \ref{tab:eda_market}), khu vực Châu Á -- Thái Bình Dương (APAC) là thị trường lớn nhất của tập đoàn với 11.002 dòng giao dịch (5.437 đơn hàng), đóng góp 3.585.744,13 USD doanh thu (chiếm 28,36\% toàn cầu) và 436.000,05 USD lợi nhuận (biên lợi nhuận 12,16\%). Đứng thứ hai là thị trường Châu Âu (EU) với 10.000 dòng giao dịch, mang lại 2.938.089,06 USD doanh thu và 372.829,74 USD lợi nhuận, đạt biên lợi nhuận lành mạnh 12,69\% nhờ duy trì mức chiết khấu trung bình thấp (10,31\%). Thị trường Hoa Kỳ (US) và Mỹ Latinh (LATAM) lần lượt đóng góp 2.297.200,86 USD và 2.164.605,17 USD doanh thu. Điểm bất thường nghiêm trọng nhất nằm ở khu vực Trung Đông và Đông Âu (EMEA): mặc dù tạo ra 806.161,31 USD doanh thu từ 5.029 dòng giao dịch, lợi nhuận thu về chỉ vỏn vẹn 43.897,97 USD (biên lợi nhuận tụt xuống mức báo động 5,45\%) do mức chiết khấu trung bình tại khu vực này bị đẩy lên tới 19,61\%. Ngược lại, thị trường Canada tuy có quy mô nhỏ (384 dòng giao dịch, doanh thu 66.928,17 USD) nhưng không áp dụng chiết khấu (0,00\%) nên đạt biên lợi nhuận cao nhất toàn hệ thống là 26,62\%.

\begin{table}[H]
\centering
\caption{Thống kê kết quả kinh doanh theo khu vực thị trường giai đoạn 2011--2014}
\label{tab:eda_market}
\small
\renewcommand{\arraystretch}{1.25}
\begin{tabularx}{\textwidth}{|l|r|r|r|r|r|r|}
\hline
\textbf{Thị trường (Market)} & \textbf{Số đơn hàng} & \textbf{Số dòng (Rows)} & \textbf{Doanh thu (USD)} & \textbf{Lợi nhuận (USD)} & \textbf{Chiết khấu BQ} & \textbf{Biên LN} \\
\hline
APAC & 5.437 & 11.002 & 3.585.744,13 & 436.000,05 & 14,88\% & 12,16\% \\
\hline
EU & 4.593 & 10.000 & 2.938.089,06 & 372.829,74 & 10,31\% & 12,69\% \\
\hline
US & 5.009 & 9.994 & 2.297.200,86 & 286.397,02 & 15,62\% & 12,47\% \\
\hline
LATAM & 5.138 & 10.294 & 2.164.605,17 & 221.643,49 & 13,55\% & 10,24\% \\
\hline
EMEA & 2.462 & 5.029 & 806.161,31 & 43.897,97 & 19,61\% & 5,45\% \\
\hline
Africa & 2.232 & 4.587 & 783.773,21 & 88.871,63 & 15,67\% & 11,34\% \\
\hline
Canada & 201 & 384 & 66.928,17 & 17.817,39 & 0,00\% & 26,62\% \\
\hline
\textbf{Tổng cộng / Toàn cục} & \textbf{25.035} & \textbf{51.290} & \textbf{12.642.501,91} & \textbf{1.467.457,29} & \textbf{14,29\%} & \textbf{11,61\%} \\
\hline
\end{tabularx}
\end{table}

Xét theo phân khúc khách hàng (\texttt{Segment}), nhóm khách hàng Cá nhân (\texttt{Consumer}) chiếm tỷ trọng áp đảo với 818 khách hàng độc lập, thực hiện 13.104 đơn hàng, mang lại 6.507.949,42 USD doanh thu (51,48\% tổng doanh thu) và 749.239,78 USD lợi nhuận (biên lợi nhuận 11,51\%). Phân khúc Doanh nghiệp (\texttt{Corporate}) gồm 476 khách hàng với 7.673 đơn hàng, đóng góp 3.824.697,52 USD doanh thu (30,25\%) và 441.208,33 USD lợi nhuận (biên lợi nhuận 11,54\%). Phân khúc Văn phòng tại gia (\texttt{Home Office}) gồm 296 khách hàng với 4.687 đơn hàng, tạo ra 2.309.854,97 USD doanh thu (18,27\%) và 277.009,18 USD lợi nhuận (đạt biên lợi nhuận cao nhất trong ba phân khúc là 11,99\%).

Xét theo cơ cấu danh mục sản phẩm, ngành hàng Công nghệ (\texttt{Technology}) dẫn đầu về cả doanh thu lẫn khả năng sinh lời với 4.744.557,50 USD doanh thu và 663.778,73 USD lợi nhuận (biên lợi nhuận 13,99\%), nổi bật nhờ hai tiểu ngành hàng chủ lực là Điện thoại (\texttt{Phones} đạt 1.706.824,14 USD doanh thu, 216.717,01 USD lợi nhuận) và Máy sao chụp (\texttt{Copiers} đạt 1.509.436,27 USD doanh thu, 258.567,55 USD lợi nhuận, biên lợi nhuận 17,13\%). Ngành hàng Đồ dùng văn phòng (\texttt{Office Supplies}) tuy có giá trị đơn vị thấp nhưng tần suất mua lại rất lớn (chiếm 31.273 dòng giao dịch), tạo ra 3.787.070,23 USD doanh thu và 518.473,83 USD lợi nhuận (biên lợi nhuận 13,69\%), trong đó nhóm Giấy in (\texttt{Paper}) và Nhãn dán (\texttt{Labels}) đạt biên lợi nhuận lần lượt tới 24,24\% và 20,54\%. Ngược lại, ngành hàng Nội thất (\texttt{Furniture}) tuy đạt doanh thu 4.110.874,19 USD nhưng lợi nhuận chỉ đạt 285.204,72 USD (biên lợi nhuận 6,94\%). Nguyên nhân trực tiếp kéo tụt hiệu quả của ngành Nội thất là tiểu ngành hàng Bàn làm việc (\texttt{Tables}): với 861 dòng giao dịch, nhóm \texttt{Tables} đạt doanh số 757.041,92 USD nhưng gây lỗ ròng tới $-64.083,39$ USD (biên lợi nhuận âm $-8,46\%$) do mức chiết khấu bình quân lên đến 29,07\% kết hợp cùng chi phí vận chuyển hàng cồng kềnh rất cao.

\subsection{Đánh giá chất lượng và rủi ro dữ liệu}

Thông qua quá trình kiểm tra chất lượng dữ liệu nguồn (Data Profiling), nhóm thực hiện phát hiện năm vấn đề kỹ thuật và rủi ro dữ liệu cần phải được xử lý triệt để trong quy trình ETL trước khi đưa vào mô hình phân tích đa chiều:

Vấn đề thứ nhất là sự bất đồng nhất nghiêm trọng về định dạng ngày tháng tại hai trường \texttt{Order Date} và \texttt{Ship Date}. Trong cùng một tệp CSV, các bản ghi từ ngày 01 đến ngày 12 hàng tháng thường được lưu theo dạng \texttt{D/M/YYYY} sử dụng dấu gạch chéo (ví dụ: \texttt{1/1/2011}), trong khi các bản ghi từ ngày 13 trở đi lại được lưu theo định dạng \texttt{DD-MM-YYYY} sử dụng dấu gạch ngang (ví dụ: \texttt{13-01-2011}). Nếu sử dụng bộ phân tích ngày tháng mặc định của hệ quản trị cơ sở dữ liệu mà không phân nhánh biểu thức chính quy (Regular Expression), quá trình chuyển đổi kiểu dữ liệu sẽ phát sinh lỗi \texttt{DateTimeParseException} hoặc đảo ngược ngày với tháng làm sai lệch hoàn toàn tính mùa vụ của chuỗi thời gian.

Vấn đề thứ hai là tỷ lệ khuyết thiếu dữ liệu (Missing Values) rất lớn tại thuộc tính \texttt{Postal Code}. Trong tổng số 51.290 dòng giao dịch, chỉ có 9.994 dòng thuộc lãnh thổ Hoa Kỳ là có mã bưu chính hợp lệ (tương ứng với 632 tổ hợp mã bưu chính duy nhất được trích xuất tại tầng Bronze), trong khi 41.296 dòng còn lại (chiếm 80,51\% toàn tập dữ liệu) mang giá trị rỗng (\texttt{NULL} hoặc chuỗi trống). Nếu sử dụng trực tiếp cột \texttt{Postal Code} chứa giá trị \texttt{NULL} làm một phần của khóa tự nhiên tổ hợp trong bảng chiều địa lý (\texttt{dim\_geography}), phép toán \texttt{UNIQUE} và phép kết nối \texttt{JOIN} chuẩn SQL sẽ thất bại do trong đại số quan hệ, biểu thức \texttt{NULL = NULL} không trả về giá trị đúng (\texttt{UNKNOWN}). Do đó, việc chuẩn hóa các giá trị khuyết thiếu về một nhãn mặc định tường minh (\texttt{'Unknown'}) là yêu cầu bắt buộc.

Vấn đề thứ ba là hiện tượng đa trị và phụ thuộc hàm không chặt chẽ trong phân cấp địa lý và danh mục sản phẩm. Cụ thể, một số thành phố có cùng tên gọi và cùng quốc gia nhưng lại xuất hiện ở các khu vực quản lý (\texttt{Market} / \texttt{Region}) khác nhau trong dữ liệu lịch sử (ví dụ các đơn hàng tại Vienna, Áo vừa xuất hiện dưới thị trường \texttt{EU} vừa xuất hiện dưới thị trường \texttt{EMEA}), hoặc cùng một mã định danh sản phẩm tự nhiên (\texttt{Product ID}) nhưng có sự khác biệt nhỏ về chuỗi tên sản phẩm (\texttt{Product Name}) giữa các năm. Nếu thiết kế ràng buộc duy nhất không đầy đủ hoặc không sử dụng kỹ thuật chọn bản ghi đại diện mới nhất (\texttt{SELECT DISTINCT ON ... ORDER BY order\_date DESC}), lệnh nạp chiều bằng \texttt{ON CONFLICT DO UPDATE} trong PostgreSQL sẽ lập tức dừng thực thi với lỗi \texttt{CardinalityViolation: ON CONFLICT DO UPDATE command cannot affect row a second time}.

Vấn đề thứ tư là rủi ro trùng lặp lô nạp (Duplicate Batch Ingestion) và ép kiểu số học. Tất cả 24 cột khi đọc từ tệp văn bản CSV vào tầng thô đều ở dạng chuỗi ký tự (\texttt{TEXT}) có thể kèm theo khoảng trắng thừa ở hai đầu. Nếu người vận hành vô tình kích hoạt quy trình nạp nhiều lần trên cùng một tệp nguồn mà không có cơ chế kiểm tra chữ ký số của tệp và cơ chế chống ghi đè tại khóa chính \texttt{row\_id}, các chỉ số cộng được như doanh thu và lợi nhuận sẽ bị nhân bội sai thực tế.

\subsection{Xác định yêu cầu và câu hỏi nghiệp vụ}

Từ kết quả phân tích khám phá và các vấn đề thực tiễn nêu trên, hệ thống Kho dữ liệu và Hỗ trợ quyết định được đặt ra yêu cầu phải chuẩn hóa 13 chỉ số đo lường hiệu quả cốt lõi (KPIs) thuộc ba nhóm tài chính, sản phẩm và khách hàng, được định nghĩa toán học chặt chẽ trong Bảng \ref{tab:kpi_definitions}.

\begin{table}[H]
\centering
\caption{Hệ thống các chỉ số đo lường hiệu quả kinh doanh (KPIs) chuẩn hóa trong hệ thống}
\label{tab:kpi_definitions}
\footnotesize
\renewcommand{\arraystretch}{1.3}
\begin{tabularx}{\textwidth}{|l|l|l|X|}
\hline
\textbf{Nhóm KPI} & \textbf{Tên chỉ số} & \textbf{Công thức xác định} & \textbf{Phân loại \& Ý nghĩa quản trị} \\
\hline
\multirow{6}{*}{Tài chính} 
 & Tổng doanh thu ($S$) & $S = \sum \text{sales}$ & Độ đo cộng được (Additive); phản ánh tổng quy mô doanh số thuần. \\ \cline{2-4}
 & Tổng lợi nhuận ($P$) & $P = \sum \text{profit}$ & Độ đo cộng được; đo lường hiệu quả sinh lời thực tế sau chi phí. \\ \cline{2-4}
 & Biên lợi nhuận ($M$) & $M = \frac{\sum \text{profit}}{\sum \text{sales}} \times 100\%$ & Độ đo không cộng được (Non-additive); đánh giá tỷ suất sinh lời trên doanh thu. \\ \cline{2-4}
 & Chiết khấu bình quân ($\bar{D}$) & $\bar{D} = \frac{1}{N}\sum \text{discount} \times 100\%$ & Độ đo bán cộng (Semi-additive); giám sát mức độ giảm giá trung bình. \\ \cline{2-4}
 & Giá trị đơn trung bình ($AOV$) & $AOV = \frac{\sum \text{sales}}{\text{CountDistinct}(\text{order\_id})}$ & Non-additive; đo lường quy mô chi tiêu bình quân trên mỗi giỏ hàng ($504,99$ USD). \\ \cline{2-4}
 & Tăng trưởng tháng ($g_{\text{MoM}}$) & $g_{\text{MoM}} = \frac{S_t - S_{t-1}}{S_{t-1}} \times 100\%$ & Non-additive; theo dõi động lượng tăng trưởng doanh thu giữa hai tháng liên tiếp. \\
\hline
\multirow{3}{*}{Sản phẩm} 
 & Sản lượng tiêu thụ ($Q$) & $Q = \sum \text{quantity}$ & Additive; tổng số đơn vị hàng hóa xuất kho (178.312 đơn vị). \\ \cline{2-4}
 & Tỷ trọng tích lũy Pareto & $P_{\text{cum}}(k) = \frac{\sum_{i=1}^k S_i}{\sum_{i=1}^M S_i} \times 100\%$ & Non-additive; xác định nhóm 20\% danh mục mang lại 80\% doanh thu/lợi nhuận. \\ \cline{2-4}
 & Độ nâng kết hợp ($\text{Lift}$) & $\text{Lift}(A \rightarrow B) = \frac{P(A \cap B)}{P(A) \cdot P(B)}$ & Đánh giá mức độ phụ thuộc mua kèm giữa hai sản phẩm để đề xuất bán chéo. \\
\hline
\multirow{3}{*}{Khách hàng} 
 & Độ mới giao dịch (\texttt{Recency}) & $R = \text{Date}_{\text{ref}} - \max(\text{order\_date})$ & Số ngày kể từ lần mua hàng gần nhất của khách hàng tính đến mốc tham chiếu. \\ \cline{2-4}
 & Tần suất mua (\texttt{Frequency}) & $F = \text{CountDistinct}(\text{order\_id})$ & Tổng số đơn đặt hàng độc lập mà một khách hàng đã thực hiện. \\ \cline{2-4}
 & Giá trị chi tiêu (\texttt{Monetary}) & $M_{\text{cust}} = \sum \text{sales}$ & Tổng giá trị doanh thu mà khách hàng đóng góp trong toàn bộ vòng đời. \\
\hline
\end{tabularx}
\end{table}

Dựa trên hệ thống KPI này, kho dữ liệu và hệ hỗ trợ quyết định được thiết kế để trả lời trực tiếp bốn nhóm câu hỏi nghiệp vụ trọng tâm. Nhóm câu hỏi thứ nhất về hiệu quả doanh thu và địa lý: Xu hướng doanh thu và lợi nhuận biến thiên như thế nào theo từng tháng, quý, năm tại từng thị trường và quốc gia, đồng thời khu vực nào đang rơi vào trạng thái doanh số cao nhưng biên lợi nhuận âm do lạm dụng chiết khấu? Nhóm câu hỏi thứ hai về quản trị danh mục sản phẩm: Những tiểu ngành hàng và mã sản phẩm SKU cụ thể nào nằm trong nhóm đứng đầu về lợi nhuận (Top Performers) hoặc nhóm gây thất thoát dòng tiền lớn nhất (Loss Makers) cần phải cắt giảm chiết khấu? Nhóm câu hỏi thứ ba về hành vi khách hàng và giỏ hàng: Cơ cấu đóng góp doanh thu theo các phân khúc Consumer, Corporate, Home Office ra sao và những cặp sản phẩm/tiểu ngành hàng nào có xác suất mua kèm cao nhất để thiết kế gói khuyến mãi? Nhóm câu hỏi thứ tư về dự báo và mô phỏng chính sách (DSS): Doanh số dự kiến trong 6 tháng tiếp theo sẽ đạt mức bao nhiêu trong khoảng tin cậy 95\%, và nếu ban giám đốc điều chỉnh tăng hoặc giảm tỷ lệ chiết khấu một mức $\Delta d$ tại một thị trường/ngành hàng cụ thể thì lợi nhuận ròng sẽ thay đổi như thế nào để ra quyết định phê duyệt hay từ chối chính sách?

\newpage
\section{Thiết kế và xây dựng Data Warehouse}

\subsection{Kiến trúc tổng thể}

Hệ thống Kho dữ liệu Global Superstore được thiết kế theo kiến trúc hồ dữ liệu đa tầng Medallion (Bronze -- Silver -- Gold) kết hợp cùng tầng phân tích hỗ trợ quyết định (\texttt{dss}), được triển khai đồng bộ trên hệ quản trị cơ sở dữ liệu quan hệ PostgreSQL 17 đóng gói trong môi trường ảo hóa Docker (\texttt{docker-compose.yml}). Thay vì nạp trực tiếp dữ liệu từ tệp văn bản phẳng vào các bảng báo cáo, kiến trúc phân tầng chia nhỏ vòng đời dữ liệu thành bốn lược đồ (schemas) độc lập về mặt logic trong cùng một cơ sở dữ liệu \texttt{global\_superstore}, giúp cô lập lỗi, bảo toàn lịch sử dữ liệu gốc và tối ưu hóa cho từng mục đích truy cập khác nhau. Sơ đồ luồng dữ liệu tổng thể của hệ thống được minh họa trong Hình \ref{fig:dw_architecture}.

\begin{figure}[H]
\centering
\begin{tikzpicture}[
    node distance=0.8cm and 0.9cm,
    box/.style={rectangle, draw=black, thick, rounded corners=3pt, minimum width=2.7cm, minimum height=1.3cm, align=center, font=\footnotesize, fill=gray!10},
    arrow/.style={-{Stealth[scale=1.0]}, thick}
]
    \node[box, fill=blue!8] (src) {\textbf{Nguồn dữ liệu}\\CSV (51.290 dòng)\\Batch Q1/2015};
    \node[box, fill=orange!12, right=of src] (bronze) {\textbf{Tầng Bronze}\\\texttt{raw\_orders}\\\texttt{ingestion} (SHA-256)};
    \node[box, fill=gray!20, right=of bronze] (silver) {\textbf{Tầng Silver}\\\texttt{silver.orders}\\(Làm sạch \& Chuẩn hóa)};
    \node[box, fill=yellow!18, right=of silver] (gold) {\textbf{Tầng Gold}\\Star Schema (1 Fact, 5 Dims)\\3 Materialized Views};
    \node[box, fill=green!12, below=0.7cm of gold] (dss) {\textbf{Tầng DSS \& BI}\\Schema \texttt{dss}, OLAP,\\Streamlit, Power BI};

    \draw[arrow] (src) -- (bronze);
    \draw[arrow] (bronze) -- (silver);
    \draw[arrow] (silver) -- (gold);
    \draw[arrow] (gold) -- (dss);
\end{tikzpicture}
\caption{Kiến trúc phân tầng Medallion của hệ thống Kho dữ liệu và DSS Global Superstore (Gợi ý: Có thể thay thế bằng sơ đồ kiến trúc chi tiết vẽ trên Draw.io nếu cần trình bày đồ họa)}
\label{fig:dw_architecture}
\end{figure}

Trong kiến trúc này, tầng \texttt{bronze} đóng vai trò là vùng tiếp nhận thô (Landing/Staging Zone), lưu trữ nguyên trạng toàn bộ các dòng bản ghi đọc từ tệp CSV dưới kiểu dữ liệu chuỗi ký tự (\texttt{TEXT}) kèm theo bảng kiểm toán lịch sử nạp \texttt{bronze.ingestion}. Tầng \texttt{silver} đóng vai trò là vùng dữ liệu đã làm sạch và chuẩn hóa (Conformed \& Cleansed Zone), nơi dữ liệu được ép kiểu chặt chẽ, chuẩn hóa định dạng ngày tháng theo chuẩn ISO, xử lý giá trị khuyết thiếu và kiểm tra các ràng buộc miền giá trị nghiệp vụ tại bảng \texttt{silver.orders}. Tầng \texttt{gold} là tầng trình diễn đa chiều (Dimensional Presentation Zone), tổ chức dữ liệu theo mô hình hình sao (Star Schema) phục vụ trực tiếp cho các truy vấn OLAP và hệ thống ba bảng tổng hợp vật lý (\texttt{mv\_sales\_mart}, \texttt{mv\_product\_mart}, \texttt{mv\_customer\_mart}). Cuối cùng, lược đồ \texttt{dss} lưu trữ bền vững siêu dữ liệu huấn luyện và kết quả đầu ra của các mô hình dự báo chuỗi thời gian cũng như khai phá luật kết hợp giỏ hàng, cung cấp dữ liệu cho dịch vụ Backend FastAPI, giao diện Streamlit và báo cáo Power BI Desktop.

\subsection{Thiết kế mô hình dữ liệu (Star Schema)}

Mô hình dữ liệu đa chiều tại tầng \texttt{gold} được thiết kế tuân thủ quy trình bốn bước của Ralph Kimball \cite{kimball2013}. Ở bước thứ nhất, quy trình nghiệp vụ được lựa chọn là hoạt động bán hàng và giao nhận đơn hàng thương mại điện tử. Ở bước thứ hai, độ mịn (Grain) của bảng sự kiện được tuyên bố ở cấp độ nguyên tử nhất: mỗi bản ghi trong bảng sự kiện đại diện cho đúng một dòng mặt hàng thuộc một đơn đặt hàng (\texttt{row\_id}). Việc giữ nguyên độ mịn ở cấp độ dòng đơn hàng thay vì gộp sớm theo tháng hay theo đơn hàng cho phép hệ thống hỗ trợ mọi truy vấn khoan sâu tới từng mã sản phẩm SKU hoặc phân tích các mặt hàng đồng xuất hiện trong cùng một mã đơn hàng (\texttt{order\_id}). Ở bước thứ ba và thứ tư, nhóm xác định năm bảng chiều (Dimensions) bao quanh một bảng sự kiện trung tâm (\texttt{gold.fact\_sales}), được minh họa trong Hình \ref{fig:star_schema} và mô tả chi tiết trong Bảng \ref{tab:star_schema_tables}.

\begin{figure}[H]
\centering
\begin{tikzpicture}[
    fact/.style={rectangle, draw=black, very thick, fill=yellow!15, minimum width=4.4cm, minimum height=2.6cm, align=left, font=\scriptsize},
    dim/.style={rectangle, draw=black, thick, fill=blue!8, minimum width=3.8cm, minimum height=1.8cm, align=left, font=\scriptsize},
    rel/.style={-{Stealth[scale=0.9]}, thick}
]
    \node[fact] (fact) at (0,0) {
        \textbf{\texttt{gold.fact\_sales} (51.290 rows)}\\
        \hline
        \texttt{PK sales\_key (BIGSERIAL)}\\
        \texttt{FK order\_date\_key, ship\_date\_key}\\
        \texttt{FK customer\_key, product\_key}\\
        \texttt{FK geography\_key, ship\_mode\_key}\\
        \texttt{DD order\_id, row\_id, order\_priority}\\
        \texttt{M: sales, quantity, discount,}\\
        \texttt{   profit, shipping\_cost, shipping\_days}
    };

    \node[dim] (date) at (0, 3.1) {
        \textbf{\texttt{gold.dim\_date} (2.922 rows)}\\
        \hline
        \texttt{PK date\_key (YYYYMMDD)}\\
        \texttt{full\_date, year, quarter,}\\
        \texttt{month, week\_of\_year, is\_weekend}
    };

    \node[dim] (cust) at (-5.3, 1.3) {
        \textbf{\texttt{gold.dim\_customer} (1.590 rows)}\\
        \hline
        \texttt{PK customer\_key (SERIAL)}\\
        \texttt{NK customer\_id (UNIQUE)}\\
        \texttt{customer\_name, segment}
    };

    \node[dim] (prod) at (5.3, 1.3) {
        \textbf{\texttt{gold.dim\_product} (10.292 rows)}\\
        \hline
        \texttt{PK product\_key (SERIAL)}\\
        \texttt{NK product\_id (UNIQUE)}\\
        \texttt{product\_name, category, sub\_category}
    };

    \node[dim] (geo) at (-5.3, -1.6) {
        \textbf{\texttt{gold.dim\_geography} (3.847 rows)}\\
        \hline
        \texttt{PK geography\_key (SERIAL)}\\
        \texttt{country, state, city, postal\_code}\\
        \texttt{market, region, regional\_manager}
    };

    \node[dim] (ship) at (5.3, -1.6) {
        \textbf{\texttt{gold.dim\_ship\_mode} (4 rows)}\\
        \hline
        \texttt{PK ship\_mode\_key (SERIAL)}\\
        \texttt{NK ship\_mode\_name (UNIQUE)}
    };

    \draw[rel] (date.south) -- node[right, font=\tiny] {1:N (Order/Ship)} (fact.north);
    \draw[rel] (cust.east) -- node[above, font=\tiny] {1:N} (fact.north west);
    \draw[rel] (prod.west) -- node[above, font=\tiny] {1:N} (fact.north east);
    \draw[rel] (geo.east) -- node[below, font=\tiny] {1:N} (fact.south west);
    \draw[rel] (ship.west) -- node[below, font=\tiny] {1:N} (fact.south east);
\end{tikzpicture}
\caption{Lược đồ hình sao (Star Schema) tại tầng Gold của Kho dữ liệu Global Superstore}
\label{fig:star_schema}
\end{figure}

Một điểm thiết kế quan trọng trong mô hình là việc sử dụng toàn bộ khóa thay thế kiểu số nguyên tự tăng (\texttt{SERIAL} / \texttt{BIGSERIAL}) làm khóa chính cho các bảng chiều và khóa ngoại trong bảng sự kiện thay vì sử dụng khóa tự nhiên dạng chuỗi ký tự từ hệ thống nguồn. Cụ thể, việc liên kết qua các cột số nguyên 4 byte như \texttt{customer\_key}, \texttt{product\_key}, \texttt{geography\_key} và \texttt{date\_key} giúp giảm đáng kể kích thước lưu trữ của bảng sự kiện \texttt{gold.fact\_sales} và tăng tốc độ thực thi các phép kết nối (\texttt{JOIN}) có đánh chỉ mục B-Tree trên PostgreSQL. Bên cạnh đó, bảng \texttt{gold.dim\_date} được thiết kế theo mẫu Chiều đóng nhiều vai trò (Role-Playing Dimension): chỉ với một bảng lịch duy nhất chứa 2.922 ngày liên tục từ năm 2011 đến năm 2018 (mã hóa khóa chính theo định dạng số nguyên thông minh \texttt{YYYYMMDD}, ví dụ \texttt{20110115}), bảng sự kiện \texttt{fact\_sales} tham chiếu đồng thời qua hai khóa ngoại là \texttt{order\_date\_key} và \texttt{ship\_date\_key}, cho phép phân tích song song theo cả thời điểm đặt hàng lẫn thời điểm giao hàng.

\begin{table}[H]
\centering
\caption{Đặc tả chi tiết các bảng trong mô hình Star Schema tại lược đồ \texttt{gold}}
\label{tab:star_schema_tables}
\footnotesize
\renewcommand{\arraystretch}{1.22}
\begin{tabularx}{\textwidth}{|l|l|r|X|}
\hline
\textbf{Tên bảng} & \textbf{Vai trò} & \textbf{Số bản ghi} & \textbf{Khóa chính (PK), Khóa tự nhiên (NK) và Các thuộc tính chính} \\
\hline
\texttt{gold.fact\_sales} & Fact Table & 51.290 & \textbf{PK:} \texttt{sales\_key}; \textbf{FK:} \texttt{order\_date\_key}, \texttt{ship\_date\_key}, \texttt{customer\_key}, \texttt{product\_key}, \texttt{geography\_key}, \texttt{ship\_mode\_key}; \textbf{DD:} \texttt{order\_id}, \texttt{row\_id}, \texttt{order\_priority}, \texttt{is\_returned}; \textbf{Measures:} \texttt{sales}, \texttt{quantity}, \texttt{discount}, \texttt{profit}, \texttt{shipping\_cost}, \texttt{shipping\_days}. \\
\hline
\texttt{gold.dim\_date} & Role-Playing Dim & 2.922 & \textbf{PK:} \texttt{date\_key} (\texttt{INTEGER YYYYMMDD}); \textbf{Thuộc tính:} \texttt{full\_date}, \texttt{day\_of\_week}, \texttt{day\_name}, \texttt{day\_of\_month}, \texttt{day\_of\_year}, \texttt{week\_of\_year}, \texttt{month}, \texttt{month\_name}, \texttt{quarter}, \texttt{year}, \texttt{year\_month}, \texttt{is\_weekend}. \\
\hline
\texttt{gold.dim\_customer} & Dimension (SCD1) & 1.590 & \textbf{PK:} \texttt{customer\_key} (\texttt{SERIAL}); \textbf{NK:} \texttt{customer\_id} (\texttt{UNIQUE}); \textbf{Thuộc tính:} \texttt{customer\_name}, \texttt{segment} (Consumer, Corporate, Home Office), \texttt{created\_at}. \\
\hline
\texttt{gold.dim\_product} & Dimension (SCD1) & 10.292 & \textbf{PK:} \texttt{product\_key} (\texttt{SERIAL}); \textbf{NK:} \texttt{product\_id} (\texttt{UNIQUE}); \textbf{Thuộc tính:} \texttt{product\_name}, \texttt{category} (3 ngành hàng), \texttt{sub\_category} (17 tiểu ngành hàng), \texttt{created\_at}. \\
\hline
\texttt{gold.dim\_geography} & Dimension (SCD1) & 3.847 & \textbf{PK:} \texttt{geography\_key} (\texttt{SERIAL}); \textbf{NK tổ hợp:} \texttt{(country, state, city, postal\_code, market, region)}; \textbf{Thuộc tính:} \texttt{regional\_manager}, \texttt{created\_at}. \\
\hline
\texttt{gold.dim\_ship\_mode} & Dimension & 4 & \textbf{PK:} \texttt{ship\_mode\_key} (\texttt{SERIAL}); \textbf{NK:} \texttt{ship\_mode\_name} (Standard Class, Second Class, First Class, Same Day). \\
\hline
\end{tabularx}
\end{table}

Các thuộc tính cấp độ giao dịch không có thuộc tính mô tả đi kèm như mã đơn hàng (\texttt{order\_id}), mã dòng (\texttt{row\_id}), độ ưu tiên (\texttt{order\_priority}) và cờ hoàn trả (\texttt{is\_returned}) được lưu trữ trực tiếp trong bảng sự kiện dưới dạng Chiều suy biến (Degenerate Dimensions). Cách thiết kế này giúp tránh việc phải tạo thêm các bảng chiều phụ có kích thước lớn tương đương bảng sự kiện mà vẫn phục vụ hiệu quả việc gom nhóm giỏ hàng theo \texttt{order\_id}. Ngoài năm độ đo tài chính gốc, bảng sự kiện còn lưu trữ độ đo dẫn xuất \texttt{shipping\_days} được tính bằng hiệu số ngày giao hàng và ngày đặt hàng (\texttt{ship\_date - order\_date}) để phục vụ phân tích hiệu suất chuỗi cung ứng.

\subsection{Quy trình ETL}

Toàn bộ quy trình ETL được điều phối tự động thông qua luồng chạy tổng thể \texttt{src/pipeline/flows/bronze\_to\_gold.py}, kết nối tuần tự ba giai đoạn xử lý từ tệp nguồn CSV đến các bảng chiều và bảng sự kiện tầng Gold.

\subsubsection{Bronze Layer}

Quy trình tại tầng Bronze được cài đặt trong mô-đun \texttt{src/etl/bronze/data\_to\_bronze.py}, kết hợp cùng hai mô-đun con là \texttt{extract.py} và \texttt{load.py}. Khi nhận đường dẫn tệp dữ liệu đầu vào, hàm \texttt{read\_raw\_data()} tiến hành đọc tệp CSV (tự động xử lý bảng mã ký tự) và gọi hàm \texttt{create\_content\_hash()} để tính giá trị băm mật mã SHA-256 độ dài 64 ký tự thập lục phân đại diện cho toàn bộ nội dung tệp (ví dụ đối với tệp gốc 51.290 dòng, mã băm thu được là \texttt{3eab4d2ae1aa2a6d...}). Ngay sau đó, hàm \texttt{is\_file\_extracted()} thực thi truy vấn \texttt{bronze/validation/check\_duplicates.sql} để kiểm tra xem mã băm \texttt{content\_hash} này đã tồn tại trong bảng kiểm toán \texttt{bronze.ingestion} hay chưa. Nếu mã băm đã tồn tại, hệ thống lập tức dừng bước nạp thô và xuất thông báo bỏ qua, đảm bảo tính lũy đẳng tuyệt đối ở cấp độ tệp tin.

Trong trường hợp tệp tin mới chưa từng được nạp, hàm \texttt{valid\_data()} kiểm tra sự hiện diện của toàn bộ 20 cột bắt buộc và bổ sung giá trị \texttt{None} cho các cột tùy chọn nếu bị khuyết. Tiếp theo, hệ thống trích xuất danh sách 149 cặp quốc gia -- thị trường duy nhất lưu vào bảng \texttt{bronze.market} và 632 tổ hợp mã bưu chính hợp lệ lưu vào bảng \texttt{bronze.postal\_code} với cơ chế \texttt{ON CONFLICT DO NOTHING}. Một bản ghi kiểm toán mới được chèn vào bảng \texttt{bronze.ingestion} lưu trữ tên tệp, mã băm SHA-256, thời gian nạp và đối tượng siêu dữ liệu kiểu \texttt{JSONB} (\texttt{\{"rows": 51290, "pin\_count": 632, "market\_count": 149\}}), trả về khóa tự tăng \texttt{ingestion\_id}. Cuối cùng, toàn bộ 51.290 bản ghi giao dịch được nạp theo lô (\texttt{chunk\_size = 100}) vào bảng \texttt{bronze.raw\_orders} kèm theo khóa ngoại \texttt{ingestion} và cột cờ trạng thái \texttt{transform} mang giá trị mặc định là \texttt{FALSE}.

\subsubsection{Silver Layer}

Giai đoạn biến đổi sang tầng Silver được điều khiển bởi mô-đun \texttt{src/etl/silver/transform.py} thông qua việc thực thi kịch bản SQL \texttt{src/sql/silver/transform/transform\_orders.sql}. Thay vì quét lại toàn bộ bảng thô, câu lệnh chỉ lọc các bản ghi mới có cờ \texttt{WHERE r.transform = FALSE} từ bảng \texttt{bronze.raw\_orders}. Tại đây, bốn phép biến đổi làm sạch cốt lõi được thực hiện đồng thời trong một giao dịch cơ sở dữ liệu (Transaction):

Thứ nhất, hệ thống cắt bỏ khoảng trắng thừa ở hai đầu chuỗi bằng hàm \texttt{TRIM()} trên tất cả các cột văn bản. Thứ hai, để giải quyết triệt để sự hỗn loạn định dạng ngày tháng tại cột \texttt{Order Date} và \texttt{Ship Date}, câu lệnh sử dụng cấu trúc phân nhánh \texttt{CASE WHEN} kết hợp toán tử khớp biểu thức chính quy (\texttt{\~}) của PostgreSQL nhằm nhận diện chính xác từng mẫu định dạng trước khi gọi hàm \texttt{TO\_DATE()}, như thể hiện trong đoạn mã SQL dưới đây:

\begin{lstlisting}[language=SQL, caption={Xu ly da dinh dang ngay thang va lam sach du lieu tai tang Silver}]
SELECT 
    CAST(TRIM(r."Row ID") AS INTEGER) AS row_id,
    TRIM(r."Order ID") AS order_id,
    CASE 
        WHEN r."Order Date" ~ '^\d{1,2}/\d{1,2}/\d{4}$' THEN TO_DATE(r."Order Date", 'DD/MM/YYYY')
        WHEN r."Order Date" ~ '^\d{1,2}-\d{1,2}-\d{4}$' THEN TO_DATE(r."Order Date", 'DD-MM-YYYY')
        WHEN r."Order Date" ~ '^\d{4}-\d{1,2}-\d{1,2}$' THEN TO_DATE(r."Order Date", 'YYYY-MM-DD')
        ELSE TO_DATE(r."Order Date", 'DD/MM/YYYY')
    END AS order_date,
    COALESCE(NULLIF(TRIM(r."Postal Code"), ''), 'Unknown') AS postal_code,
    CAST(TRIM(r."Sales") AS NUMERIC(12, 4)) AS sales,
    CAST(TRIM(r."Quantity") AS INTEGER) AS quantity,
    CAST(TRIM(r."Discount") AS NUMERIC(5, 4)) AS discount,
    CAST(TRIM(r."Profit") AS NUMERIC(12, 4)) AS profit
FROM bronze.raw_orders r
WHERE r.transform = FALSE
ON CONFLICT (row_id) DO NOTHING;
\end{lstlisting}

Thứ ba, các giá trị rỗng hoặc chuỗi trống tại cột \texttt{Postal Code} được chuẩn hóa đồng nhất về chuỗi \texttt{'Unknown'} thông qua tổ hợp hàm \texttt{COALESCE(NULLIF(TRIM(...), ''), 'Unknown')}; tương tự, cột \texttt{Order Priority} nếu bị rỗng sẽ được gán mặc định là \texttt{'Medium'}. Thứ tư, các trường định lượng được ép kiểu tường minh sang \texttt{NUMERIC(12,4)} đối với doanh thu, lợi nhuận, chi phí vận chuyển và \texttt{NUMERIC(5,4)} đối với chiết khấu. Bảng đích \texttt{silver.orders} được bảo vệ bởi bốn ràng buộc toàn vẹn miền giá trị gồm \texttt{chk\_silver\_sales CHECK (sales >= 0)}, \texttt{chk\_silver\_quantity CHECK (quantity > 0)}, \texttt{chk\_silver\_discount CHECK (discount >= 0 AND discount <= 1.0)} và \texttt{chk\_silver\_dates CHECK (ship\_date >= order\_date)}. Ngay sau khi lệnh \texttt{INSERT ... ON CONFLICT (row\_id) DO NOTHING} hoàn tất, câu lệnh \texttt{UPDATE bronze.raw\_orders SET transform = TRUE WHERE transform = FALSE} được thực thi để đánh dấu các bản ghi đã được xử lý xong.

\subsubsection{Gold Layer}

Quá trình nạp dữ liệu vào mô hình hình sao tầng Gold được chia thành hai bước tuần tự: nạp các bảng chiều (\texttt{src/etl/gold/load\_dimensions.py}) và nạp bảng sự kiện (\texttt{src/etl/gold/load\_facts.py}). Ở bước nạp chiều, bảng lịch \texttt{gold.dim\_date} được khởi tạo tự động bằng hàm sinh chuỗi thời gian \texttt{generate\_series('2011-01-01'::DATE, '2018-12-31'::DATE, '1 day'::INTERVAL)} của PostgreSQL, trích xuất đầy đủ các thuộc tính thứ, ngày, tuần ISO, tháng, quý, năm và cờ cuối tuần cho 2.922 ngày liên tiếp. Đối với bảng \texttt{gold.dim\_customer} và \texttt{gold.dim\_product}, hệ thống áp dụng chiến lược Chiều biến đổi chậm loại 1 (SCD Type 1) thông qua cú pháp \texttt{SELECT DISTINCT ON (customer\_id) ... ORDER BY customer\_id, order\_date DESC} kết hợp mệnh đề \texttt{ON CONFLICT DO UPDATE}. Kỹ thuật này đảm bảo nếu một mã khách hàng hoặc mã sản phẩm xuất hiện nhiều lần với tên gọi hoặc phân khúc cập nhật theo thời gian, hệ thống sẽ luôn lấy bản ghi có ngày đặt hàng mới nhất để cập nhật mà không gây lỗi trùng khóa trong cùng một lô nạp.

Đối với bảng chiều địa lý \texttt{gold.dim\_geography}, để khắc phục hiện tượng một số thành phố thuộc nhiều vùng thị trường khác nhau, nhóm thiết lập ràng buộc duy nhất gồm sáu thuộc tính \texttt{UNIQUE (country, state, city, postal\_code, market, region)} và trích xuất ra 3.847 vị trí địa lý phân biệt. Ở bước nạp bảng sự kiện (\texttt{src/sql/gold/facts/load\_fact\_sales.sql}), hệ thống thực hiện kết nối (\texttt{JOIN}) bảng \texttt{silver.orders} với các bảng chiều để tra cứu toàn bộ các khóa thay thế (\texttt{customer\_key}, \texttt{product\_key}, \texttt{geography\_key}, \texttt{ship\_mode\_key}), chuyển đổi ngày đặt hàng và ngày giao hàng thành khóa số nguyên \texttt{YYYYMMDD} bằng biểu thức \texttt{CAST(TO\_CHAR(s.order\_date, 'YYYYMMDD') AS INTEGER)}, tính toán số ngày giao hàng \texttt{shipping\_days} và nạp vào \texttt{gold.fact\_sales} với mệnh đề \texttt{ON CONFLICT (row\_id) DO NOTHING}.

\subsubsection{Cơ chế nạp gia tăng (Incremental Loading) và bắt thay đổi dữ liệu}

Để chứng minh khả năng vận hành trong môi trường sản xuất thực tế khi có các đơn hàng mới phát sinh định kỳ mà không cần xóa hay nạp lại toàn bộ kho dữ liệu lịch sử, nhóm đã xây dựng mô-đun tạo lô dữ liệu mới (\texttt{src/etl/incremental/generate\_batch.py}) và kịch bản kiểm thử nạp gia tăng tự động (\texttt{src/etl/incremental/demo\_incremental.py}). Cơ chế bắt thay đổi dữ liệu (Change Data Capture -- CDC) của hệ thống dựa trên sự kết hợp giữa bảng kiểm toán \texttt{bronze.ingestion}, cột cờ đánh dấu \texttt{transform} tại bảng \texttt{bronze.raw\_orders} và cơ chế Upsert lũy đẳng tại tầng Silver và Gold.

Khi kích hoạt lệnh \texttt{python -m src.etl.incremental.demo\_incremental}, hệ thống đưa tệp dữ liệu mới \texttt{dataset/incremental\_orders\_2015\_q1.csv} gồm 500 dòng giao dịch phát sinh trong Quý 1/2015 (có \texttt{Row ID} từ 60001 đến 60500) vào đường ống. Tầng Bronze nhận diện mã băm SHA-256 mới (\texttt{975dd415be1161ee...}), ghi nhận đợt nạp thứ hai vào \texttt{bronze.ingestion} và chỉ chèn thêm 500 dòng mới vào \texttt{bronze.raw\_orders} với cờ \texttt{transform = FALSE} (trong khi 51.290 dòng cũ vẫn giữ nguyên \texttt{transform = TRUE}). Khi tầng Silver thực thi, mệnh đề \texttt{WHERE r.transform = FALSE} giúp hệ thống chỉ quét và biến đổi đúng 500 bản ghi mới này để chèn vào \texttt{silver.orders}. Tiếp đó, tầng Gold tự động cập nhật các bảng chiều, nạp bổ sung 500 dòng mới vào \texttt{gold.fact\_sales} và phát lệnh \texttt{REFRESH MATERIALIZED VIEW} để làm mới toàn bộ các Data Marts. Kết quả đối chứng thực nghiệm trước và sau khi nạp gia tăng được trình bày trong Bảng \ref{tab:incremental_results}. Ngoài ra, kịch bản còn cung cấp tham số \texttt{--reset} cho phép thu hồi chính xác lô dữ liệu thử nghiệm (\texttt{row\_id >= 60000}) và khôi phục kho dữ liệu về trạng thái chuẩn 51.290 dòng ban đầu.

\begin{table}[H]
\centering
\caption{Kết quả thực nghiệm quy trình nạp gia tăng (Incremental Loading) với lô dữ liệu Quý 1/2015}
\label{tab:incremental_results}
\small
\renewcommand{\arraystretch}{1.25}
\begin{tabularx}{\textwidth}{|X|r|r|r|}
\hline
\textbf{Chỉ số giám sát kho dữ liệu} & \textbf{Trước nạp (Baseline)} & \textbf{Sau nạp gia tăng} & \textbf{Biến thiên (Delta)} \\
\hline
Số đợt nạp trong \texttt{bronze.ingestion} & 1 đợt & 2 đợt & $+1$ đợt \\
\hline
Tổng số bản ghi trong \texttt{bronze.raw\_orders} & 51.290 dòng & 51.790 dòng & $+500$ dòng \\
\hline
Tổng số bản ghi trong \texttt{silver.orders} & 51.290 dòng & 51.790 dòng & $+500$ dòng \\
\hline
Tổng số bản ghi trong \texttt{gold.fact\_sales} & 51.290 dòng & 51.790 dòng & $+500$ dòng \\
\hline
Tổng doanh thu lũy kế (USD) & 12.642.501,91 & 13.256.279,10 & $+613.777,19$ \\
\hline
Tổng lợi nhuận ròng lũy kế (USD) & 1.467.457,29 & 1.606.297,96 & $+138.840,67$ \\
\hline
\end{tabularx}
\end{table}

\subsection{OLAP và các Data Mart / Materialized Views}

Trên nền tảng mô hình hình sao tầng Gold, nhóm đã xây dựng bộ 7 kịch bản truy vấn phân tích đa chiều OLAP chuẩn mực đặt tại thư mục \texttt{src/sql/gold/marts/} và được thực thi tự động qua trình chạy \texttt{run\_olap\_suite.py}. Mỗi truy vấn giải quyết một góc nhìn quản trị cụ thể và khai thác tối đa các toán tử tập hợp nâng cao của SQL trên PostgreSQL 17 \cite{postgresql17}, với kết quả thực nghiệm tổng hợp trong Bảng \ref{tab:olap_suite}:

\begin{table}[H]
\centering
\caption{Tổng hợp kết quả thực thi bộ 7 phép toán phân tích đa chiều OLAP trên tầng Gold}
\label{tab:olap_suite}
\footnotesize
\renewcommand{\arraystretch}{1.22}
\begin{tabularx}{\textwidth}{|l|l|r|r|X|}
\hline
\textbf{Phép toán OLAP} & \textbf{Tệp kịch bản SQL} & \textbf{Số dòng trả về} & \textbf{Thời gian} & \textbf{Kỹ thuật SQL \& Phát hiện nghiệp vụ chính} \\
\hline
1. Roll-up & \texttt{olap\_rollup.sql} & 1.829 & 49,2 ms & Dùng \texttt{GROUP BY ROLLUP (year, quarter, market, country)} tổng hợp phân cấp từ quốc gia lên tổng toàn cục (12,64 triệu USD). \\
\hline
2. Drill-down & \texttt{olap\_drilldown.sql} & 20 & 32,8 ms & Khoan sâu ngành \texttt{Technology} tới từng mã SKU; phát hiện máy in 3D \texttt{Cubify CubeX} lỗ nặng nhất ($-8.879,97$ USD do chiết khấu 53,33\%). \\
\hline
3. Slice \& Dice & \texttt{olap\_slice\_dice.sql} & 36 \& 8 & 28,5 ms & Cắt lát thị trường \texttt{APAC} năm 2014 và trích xuất khối con \texttt{(APAC, EU)} $\times$ \texttt{(Furniture, Technology)} trong Quý 3, Quý 4. \\
\hline
4. Pivot & \texttt{olap\_pivot.sql} & 17 & 19,4 ms & Xoay trục 4 Quý (\texttt{Q1--Q4}) thành 4 cột ngang bằng \texttt{SUM(CASE WHEN quarter = k...)} cho 17 tiểu ngành hàng trong năm 2014. \\
\hline
5. MoM Growth & \texttt{olap\_growth\_mom.sql} & 48 & 16,1 ms & Dùng hàm cửa sổ \texttt{LAG(current\_sales, 1) OVER (ORDER BY year, month)} tính tốc độ tăng trưởng doanh thu liên tiếp 48 tháng. \\
\hline
6. Top Ranking & \texttt{olap\_top\_products.sql} & 15 & 24,7 ms & Dùng \texttt{DENSE\_RANK() OVER (PARTITION BY category ORDER BY SUM(sales) DESC)} trích xuất Top 5 sản phẩm đầu bảng mỗi ngành. \\
\hline
7. CUBE & \texttt{olap\_cube.sql} & 78 & 21,3 ms & Dùng \texttt{GROUP BY CUBE (year, category, market)} sinh đồng thời $2^3 = 8$ cấp độ tổng hợp đa chiều chỉ trong một lần quét bảng. \\
\hline
\end{tabularx}
\end{table}

Để tăng tốc hơn nữa cho các bảng điều khiển BI thường xuyên truy vấn lặp lại các chỉ số tổng hợp, nhóm đã thiết kế và khởi tạo 3 Bảng tổng hợp vật lý (Materialized Views) trong kịch bản \texttt{src/sql/gold/marts/create\_materialized\_views.sql}. Bảng thứ nhất là \texttt{gold.mv\_sales\_mart} (lưu trữ 7.424 dòng tổng hợp trước theo Năm, Quý, Tháng, Thị trường, Vùng, Quốc gia và Phân khúc khách hàng, được đánh chỉ mục trên thời gian và địa lý). Bảng thứ hai là \texttt{gold.mv\_product\_mart} (lưu trữ 29.363 dòng tổng hợp hiệu suất bán hàng, lợi nhuận và biên lợi nhuận của từng sản phẩm theo từng thị trường và năm). Bảng thứ ba là \texttt{gold.mv\_customer\_mart} (lưu trữ đúng 1.590 dòng tương ứng với 1.590 khách hàng, tính toán sẵn ba chỉ số phân tích hành vi RFM gồm \texttt{recency\_days} tính tới mốc tham chiếu \texttt{2015-01-01}, tần suất đặt hàng \texttt{frequency}, tổng chi tiêu \texttt{monetary} và giá trị đơn trung bình \texttt{avg\_order\_value}). Nhờ việc lưu trữ vật lý kết quả tính toán sẵn trên đĩa cứng kèm chỉ mục B-Tree, các truy vấn từ giao diện đồ họa vào ba Data Mart này đạt thời gian phản hồi dưới 5 ms.

\subsection{Kiểm thử và kiểm tra tính đúng đắn của dữ liệu}

Nhằm đảm bảo dữ liệu sau khi trải qua các bước biến đổi không bị thất thoát, biến dạng hay vi phạm toàn vẹn tham chiếu, hệ thống tích hợp hai bộ kiểm thử tự động chạy ngay trong luồng pipeline là \texttt{src/etl/silver/validate.py} và \texttt{src/etl/gold/validate.py}. Tại tầng Silver, ba kịch bản kiểm tra được thực thi gồm: \texttt{check\_nulls.sql} (xác nhận 0 giá trị \texttt{NULL} trên toàn bộ 9 cột bắt buộc như \texttt{row\_id}, \texttt{order\_id}, \texttt{order\_date}, \texttt{customer\_id}, \texttt{product\_id}, \texttt{sales}, \texttt{quantity}), \texttt{check\_duplicates.sql} (nhóm theo \texttt{row\_id} với điều kiện \texttt{HAVING COUNT(*) > 1} và xác nhận 0 dòng trùng lặp) và \texttt{check\_data\_quality.sql} (kiểm tra các quy tắc nghiệp vụ, xác nhận 0 đơn hàng có ngày giao trước ngày đặt, 0 dòng có doanh thu âm, 0 dòng có số lượng $\le 0$ và 0 dòng có chiết khấu ngoài đoạn $[0, 1]$).

\begin{table}[H]
\centering
\caption{Kết quả đối soát dữ liệu (Data Reconciliation) và kiểm định chất lượng giữa các tầng}
\label{tab:data_reconciliation}
\small
\renewcommand{\arraystretch}{1.25}
\begin{tabularx}{\textwidth}{|X|r|r|r|c|}
\hline
\textbf{Tiêu chí kiểm định \& Đối soát} & \textbf{Tầng Silver (\texttt{orders})} & \textbf{Tầng Gold (\texttt{fact\_sales})} & \textbf{Độ lệch} & \textbf{Trạng thái} \\
\hline
Tổng số dòng bản ghi (\texttt{total\_rows}) & 51.290 & 51.290 & 0 & Đạt (PASS) \\
\hline
Tổng doanh thu \texttt{SUM(sales)} (USD) & 12.642.501,91 & 12.642.501,91 & 0,00 & Đạt (PASS) \\
\hline
Tổng lợi nhuận \texttt{SUM(profit)} (USD) & 1.467.457,29 & 1.467.457,29 & 0,00 & Đạt (PASS) \\
\hline
Tổng sản lượng \texttt{SUM(quantity)} (Đơn vị) & 178.312 & 178.312 & 0 & Đạt (PASS) \\
\hline
Số bản ghi vi phạm độ mịn (\texttt{Duplicate row\_id}) & 0 & 0 & 0 & Đạt (PASS) \\
\hline
Số khóa ngoại mồ côi (\texttt{Orphan Foreign Keys}) & -- & 0 & 0 & Đạt (PASS) \\
\hline
Số bản ghi vi phạm miền giá trị nghiệp vụ & 0 & 0 & 0 & Đạt (PASS) \\
\hline
\end{tabularx}
\end{table}

Tại tầng Gold, bộ kiểm định thực thi ba kịch bản đối soát nghiêm ngặt: \texttt{check\_fact\_grain.sql} xác nhận bảng \texttt{gold.fact\_sales} bảo toàn nguyên vẹn độ mịn 1 dòng trên mỗi \texttt{row\_id}; \texttt{check\_foreign\_keys.sql} kiểm tra cả 6 cột khóa ngoại (\texttt{order\_date\_key}, \texttt{ship\_date\_key}, \texttt{customer\_key}, \texttt{product\_key}, \texttt{geography\_key}, \texttt{ship\_mode\_key}) và xác nhận không có bất kỳ bản ghi mồ côi nào (\texttt{orphan\_count = 0}); cuối cùng, \texttt{check\_aggregates.sql} thực hiện phép hợp \texttt{UNION ALL} để đối chiếu trực diện tổng số dòng, tổng doanh thu, tổng lợi nhuận và tổng số lượng hàng hóa giữa \texttt{silver.orders} và \texttt{gold.fact\_sales}. Kết quả thực nghiệm trong Bảng \ref{tab:data_reconciliation} chứng minh sự khớp nhau tuyệt đối đến từng xu lẻ (sai số bằng 0,00 USD), khẳng định tính đúng đắn và độ tin cậy của toàn bộ quy trình xây dựng Kho dữ liệu.

\newpage
\section{Phát triển Decision Support System (DSS)}

\subsection{Mục tiêu và kiến trúc DSS}

Nếu như Kho dữ liệu (Data Warehouse) và các truy vấn OLAP tập trung giải quyết các câu hỏi mang tính mô tả và chẩn đoán quá khứ (\textit{Chuyện gì đã xảy ra và tại sao?}), thì Hệ hỗ trợ quyết định (Decision Support System -- DSS) được phát triển nhằm nâng cấp năng lực của hệ thống lên mức độ phân tích dự đoán (\textit{Predictive Analytics: Chuyện gì sẽ xảy ra tiếp theo?}) và phân tích đề xuất hành động (\textit{Prescriptive Analytics: Doanh nghiệp nên hành động như thế nào là tối ưu?}). Mục tiêu cốt lõi của phân hệ DSS trong đề tài là cung cấp cho nhà quản trị của Global Superstore các mô hình toán học và thuật toán khai phá dữ liệu vận hành trực tiếp trên nguồn dữ liệu sạch của tầng Gold, từ đó tự động sinh ra các dự báo định lượng, các gợi ý ghép gói sản phẩm bán chéo và các khuyến nghị điều chỉnh chính sách chiết khấu.

Về mặt phương pháp luận, kiến trúc DSS của hệ thống được xây dựng dựa trên sự kết hợp giữa mô hình cấu trúc ba thành phần cổ điển của Turban \cite{turban2014} và chu trình ra quyết định bốn giai đoạn của Herbert Simon \cite{simon1960} (Thám sát -- Thiết kế -- Lựa chọn -- Thực thi). Cụ thể, phân hệ Quản lý dữ liệu (Data Management Subsystem) bao gồm lược đồ \texttt{gold} (cung cấp dữ liệu lịch sử đầu vào) và lược đồ chuyên biệt \texttt{dss} (được khởi tạo bởi \texttt{src/sql/gold/schema/create\_dss\_schema.sql}) gồm bốn bảng cơ sở dữ liệu là \texttt{dss.forecast\_run}, \texttt{dss.forecast\_result}, \texttt{dss.association\_run} và \texttt{dss.association\_rule}. Phân hệ Quản lý mô hình (Model Management Subsystem) được đóng gói thành ba gói chương trình Python độc lập trong thư mục \texttt{src/dss/} gồm mô-đun dự báo chuỗi thời gian (\texttt{predictive/run.py}), mô-đun khai phá luật kết hợp giỏ hàng (\texttt{association/run.py}) và mô-đun mô phỏng kinh tế lượng (\texttt{what\_if/simulate.py}). Phân hệ Giao diện người dùng (User Interface Subsystem) kết nối trực tiếp với tầng mô hình và cơ sở dữ liệu thông qua dịch vụ RESTful API (\texttt{src/app/api/main.py}), ứng dụng Web tương tác bốn phân trang (\texttt{src/app/main.py}) và bảng điều khiển Power BI.

\subsection{Predictive Analytics}

\subsubsection{Bài toán}

Trong quản trị chuỗi cung ứng và kế hoạch tài chính bán lẻ, việc dự báo chính xác doanh số bán hàng và sản lượng hàng hóa tiêu thụ trong các tháng tiếp theo giữ vai trò quyết định đối với việc phân bổ ngân sách mua hàng, chuẩn bị diện tích kho bãi và tối ưu hóa dòng tiền lưu động. Bài toán đặt ra cho mô-đun Predictive Analytics là: từ chuỗi dữ liệu lịch sử 48 tháng liên tục (từ tháng 01/2011 đến tháng 12/2014) được tổng hợp từ 51.290 giao dịch trong bảng sự kiện \texttt{gold.fact\_sales} và bảng chiều thời gian \texttt{gold.dim\_date}, hãy xây dựng mô hình dự báo chuỗi thời gian để ước lượng giá trị doanh thu thuần hàng tháng (\texttt{total\_sales}) hoặc sản lượng tiêu thụ (\texttt{total\_quantity}) cho tầm nhìn $H = 6$ tháng tiếp theo (từ tháng 01/2015 đến tháng 06/2015), đồng thời xác định dải khoảng tin cậy thống kê 95\% (Lower Bound và Upper Bound) nhằm định lượng biên độ rủi ro dao động của thị trường.

\subsubsection{Phương pháp}

Quy trình dự báo được cài đặt trong mô-đun \texttt{src/dss/predictive/run.py}. Trước hết, hàm \texttt{extract\_sales\_series()} thực thi truy vấn SQL kết nối \texttt{gold.fact\_sales} với \texttt{gold.dim\_date}, gom nhóm theo \texttt{d.year}, \texttt{d.month}, \texttt{d.year\_month} và sắp xếp theo đúng trình tự thời gian để thu được chuỗi quan sát $Y = \{y_0, y_1, \dots, y_{n-1}\}$ với $n = 48$ điểm dữ liệu tháng. Để đánh giá khách quan năng lực khái quát hóa của mô hình và tránh hiện tượng rò rỉ dữ liệu tương lai (Data Leakage), chuỗi thời gian được phân tách nghiêm ngặt theo thứ tự thời gian với tỷ lệ 75\% huấn luyện và 25\% kiểm thử: tập huấn luyện (Train Set) gồm 36 tháng đầu tiên từ \texttt{2011-01-01} đến \texttt{2013-12-01}, và tập kiểm thử (Test Set) gồm 12 tháng của năm cuối cùng từ \texttt{2014-01-01} đến \texttt{2014-12-01}.

Về thuật toán dự báo, hệ thống triển khai phương pháp San bằng mũ kết hợp xu thế tuyến tính (Exponential Smoothing with Linear Trend) \cite{hyndman2021} và đối sánh với mô hình cơ sở Trung bình trượt 3 tháng (Moving Average Baseline với cửa sổ $w=3$). Với hệ số làm mượt $\alpha \in (0, 1)$ (được thiết lập $\alpha = 0,4$ qua thực nghiệm để cân bằng giữa trọng số của các quan sát gần nhất và tính ổn định dài hạn), chuỗi giá trị làm mượt $S_t$ được tính đệ quy theo công thức:
\begin{equation}
    S_0 = y_0, \quad S_t = \alpha \cdot y_t + (1 - \alpha) \cdot S_{t-1} \quad \forall t \in \{1, 2, \dots, n-1\}
\end{equation}
Hệ số góc phản ánh tốc độ tăng trưởng xu thế trung bình theo tháng ($b$) trên toàn chuỗi và giá trị dự báo điểm $\hat{y}_{n+h}$ tại bước dự báo tương lai $h \in \{1, 2, \dots, H\}$ được xác định bởi:
\begin{equation}
    b = \frac{S_{n-1} - S_0}{n}, \quad \hat{y}_{n+h} = S_{n-1} + b \cdot h
\end{equation}
Gọi $\sigma_e = \sqrt{\frac{1}{|T_{\text{test}}|}\sum_{t \in T_{\text{test}}} (y_t - \hat{y}_t - \bar{e})^2}$ là độ lệch chuẩn của phần dư dự báo trên tập kiểm thử, khoảng tin cậy 95\% (tương ứng với phân vị chuẩn $z_{0,975} = 1,96$) được xác định bởi cận dưới $LB_{n+h} = \max(0, \hat{y}_{n+h} - 1,96 \cdot \sigma_e)$ và cận trên $UB_{n+h} = \hat{y}_{n+h} + 1,96 \cdot \sigma_e$. Chất lượng mô hình trên tập kiểm thử được định lượng thông qua ba độ đo sai số chuẩn gồm Sai số tuyệt đối trung bình (MAE), Căn bậc hai sai số toàn phương trung bình (RMSE) và Sai số phần trăm tuyệt đối trung bình (MAPE):
\begin{equation}
    \text{MAE} = \frac{1}{m}\sum_{t=1}^{m} |y_t - \hat{y}_t|, \quad
    \text{RMSE} = \sqrt{\frac{1}{m}\sum_{t=1}^{m} (y_t - \hat{y}_t)^2}, \quad
    \text{MAPE} = \frac{100\%}{m}\sum_{t=1}^{m} \left|\frac{y_t - \hat{y}_t}{y_t}\right|
\end{equation}

\subsubsection{Thực nghiệm và kết quả}

Khi thực thi lệnh \texttt{python -m src.dss.predictive.run --target sales --horizon 6}, hệ thống tiến hành huấn luyện, đánh giá trên tập kiểm thử 12 tháng năm 2014 và tự động lưu thông tin siêu dữ liệu vào bảng \texttt{dss.forecast\_run} cùng 6 mốc dự báo tương lai vào bảng \texttt{dss.forecast\_result}. Trên tập kiểm thử năm 2014 (nơi doanh thu thực tế hàng tháng dao động mạnh từ 180.000 USD đến 555.279 USD), mô hình đạt chỉ số $\text{MAE} = 48.591,40\text{ USD}$, $\text{RMSE} = 59.185,84\text{ USD}$ và tỷ lệ sai số tương đối $\text{MAPE} = 14,72\%$ (tương ứng độ chính xác trung bình đạt $85,28\%$, vượt ngưỡng yêu cầu $\text{MAPE} < 15\%$ đối với dữ liệu bán lẻ đa quốc gia có biến động mùa vụ cao). Chi tiết kết quả dự báo cho 6 tháng đầu năm 2015 được trình bày trong Bảng \ref{tab:forecast_results}.

\begin{table}[H]
\centering
\caption{Kết quả dự báo doanh thu 6 tháng đầu năm 2015 lưu trong bảng \texttt{dss.forecast\_result}}
\label{tab:forecast_results}
\small
\renewcommand{\arraystretch}{1.25}
\begin{tabularx}{\textwidth}{|c|c|r|r|r|}
\hline
\textbf{Kỳ dự báo (Tháng)} & \textbf{Ngày mốc (\texttt{forecast\_date})} & \textbf{Dự báo trung tâm (USD)} & \textbf{Cận dưới 95\% CI (USD)} & \textbf{Cận trên 95\% CI (USD)} \\
\hline
Tháng 01/2015 ($h=1$) & \texttt{2015-01-01} & 491.719,91 & 379.453,66 & 603.986,17 \\
\hline
Tháng 02/2015 ($h=2$) & \texttt{2015-02-01} & 499.736,68 & 387.470,43 & 612.002,93 \\
\hline
Tháng 03/2015 ($h=3$) & \texttt{2015-03-01} & 507.753,44 & 395.487,19 & 620.019,69 \\
\hline
Tháng 04/2015 ($h=4$) & \texttt{2015-04-01} & 515.770,21 & 403.503,96 & 628.036,46 \\
\hline
Tháng 05/2015 ($h=5$) & \texttt{2015-05-01} & 523.786,97 & 411.520,72 & 636.053,22 \\
\hline
Tháng 06/2015 ($h=6$) & \texttt{2015-06-01} & 531.803,73 & 419.537,48 & 644.069,98 \\
\hline
\end{tabularx}
\end{table}

Kết quả dự báo cho thấy xu hướng doanh số toàn cầu trong nửa đầu năm 2015 tiếp tục duy trì đà tăng trưởng tích cực với mức tăng trung bình khoảng $+8.016,76\text{ USD/tháng}$, tiến dần từ mức 491.719,91 USD vào tháng 01/2015 lên 531.803,73 USD vào tháng 06/2015. Dải tin cậy 95\% cung cấp cho bộ phận kế hoạch tài chính kịch bản thận trọng (cận dưới quanh mức 379.453 -- 419.537 USD) và kịch bản khả quan (cận trên quanh mức 603.986 -- 644.069 USD) để chủ động chuẩn bị nguồn cung hàng hóa.

\subsection{Association Rule Mining}

\subsubsection{Bài toán}

Trong thương mại điện tử, việc gia tăng Giá trị đơn hàng trung bình (Average Order Value -- AOV) trên tập khách hàng hiện hữu mang lại hiệu quả chi phí cao hơn nhiều so với việc tìm kiếm khách hàng mới. Bài toán đặt ra là: từ 25.035 giỏ hàng giao dịch thực tế (mỗi giỏ hàng tương ứng với một mã \texttt{order\_id} chứa một tập hợp các mặt hàng được mua cùng lúc), hệ thống cần khai phá các luật kết hợp dạng $A \rightarrow B$ (nếu khách hàng mua mục hàng $A$ thì có xu hướng mua kèm mục hàng $B$ trong cùng đơn hàng) nhằm phục vụ hệ thống gợi ý bán chéo (Cross-selling) và thiết kế các gói sản phẩm khuyến mãi (Product Bundling).

Một thách thức kỹ thuật lớn của bộ dữ liệu Global Superstore là hiện tượng thưa dữ liệu (Data Sparsity) khi xét ở cấp độ mã sản phẩm chi tiết: với 10.292 mặt hàng khác nhau phân bổ trên 25.035 đơn hàng, xác suất hai sản phẩm cụ thể bất kỳ xuất hiện chung trong hơn 0,5\% tổng số đơn hàng (tức $\ge 125$ đơn hàng) là bằng không. Nếu áp dụng cứng nhắc một ngưỡng độ hỗ trợ phần trăm cao cho mọi cấp độ, thuật toán sẽ trả về tập luật rỗng. Do đó, bài toán đòi hỏi một chiến lược khai phá phân cấp hai tầng: vừa khai phá ở cấp độ Tiểu ngành hàng (\texttt{sub\_category}) để hoạch định chiến lược trưng bày gian hàng vĩ mô, vừa khai phá ở cấp độ Tên sản phẩm chi tiết (\texttt{product\_name}) với ngưỡng đếm tuyệt đối phù hợp để phát hiện các cặp sản phẩm cụ thể có mức độ gắn kết cao.

\subsubsection{Phương pháp}

Mô-đun khai phá giỏ hàng được hiện thực trong \texttt{src/dss/association/run.py}. Ở bước trích xuất giỏ hàng, hàm \texttt{extract\_baskets(engine, level)} thực thi truy vấn \texttt{SELECT DISTINCT f.order\_id, p.product\_name (hoặc p.sub\_category) AS item\_name FROM gold.fact\_sales f JOIN gold.dim\_product p ON f.product\_key = p.product\_key}, sau đó gom nhóm theo \texttt{order\_id} để tạo danh sách $N = 25.035$ tập giao dịch $T = \{t_1, t_2, \dots, t_N\}$. Đối với mỗi cặp mục hàng $(A, B)$ xuất hiện đồng thời trong ít nhất $c_{\min}$ giỏ hàng (\texttt{min\_count}), hệ thống tính toán ba độ đo chuẩn của thuật toán Apriori \cite{agrawal1994} gồm Độ hỗ trợ (Support), Độ tin cậy (Confidence) và Độ nâng (Lift):
\begin{equation}
    \text{Support}(A \rightarrow B) = \frac{|\{t \in T : \{A, B\} \subseteq t\}|}{N} = P(A \cap B)
\end{equation}
\begin{equation}
    \text{Confidence}(A \rightarrow B) = \frac{\text{Support}(A \rightarrow B)}{\text{Support}(A)} = P(B \mid A)
\end{equation}
\begin{equation}
    \text{Lift}(A \rightarrow B) = \frac{\text{Confidence}(A \rightarrow B)}{\text{Support}(B)} = \frac{P(A \cap B)}{P(A) \cdot P(B)}
\end{equation}
Trong đó, chỉ số $\text{Lift}(A \rightarrow B) > 1$ chứng tỏ việc mua $A$ làm tăng xác suất mua kèm $B$ cao hơn gấp $\text{Lift}$ lần so với xác suất mua $B$ ngẫu nhiên độc lập. Các luật thỏa mãn đồng thời ba ngưỡng lọc $\text{count} \ge c_{\min}$, $\text{Support} \ge s_{\min}$, $\text{Confidence} \ge c_{\text{conf}}$ và $\text{Lift} \ge l_{\min}$ được sắp xếp giảm dần theo hệ số Lift và ghi nhận vào hai bảng \texttt{dss.association\_run} và \texttt{dss.association\_rule}.

\subsubsection{Thực nghiệm và kết quả}

Nhóm tiến hành thực nghiệm trên toàn bộ 25.035 giao dịch ở cả hai cấp độ phân tích. Ở cấp độ thứ nhất -- cấp độ Tiểu ngành hàng (\texttt{--level sub\_category --min-support 0.005 --min-confidence 0.10 --min-lift 1.05}), thuật toán phát hiện được 14 luật kết hợp có ý nghĩa quản trị danh mục rõ nét, trong đó 8 luật tiêu biểu nhất được trình bày trong Bảng \ref{tab:assoc_subcat}. Kết quả chỉ ra rằng khách hàng mua Máy sao chụp (\texttt{Copiers}) có xu hướng mua kèm Nhãn dán (\texttt{Labels}, $\text{Lift} = 1,17$), Tủ sách (\texttt{Bookcases}, $\text{Lift} = 1,13$) và Phong bì (\texttt{Envelopes}, $\text{Lift} = 1,11$); trong khi khách hàng mua Thiết bị văn phòng (\texttt{Appliances}) thường mua kèm Giấy in (\texttt{Paper}, $\text{Confidence} = 15,0\%, \text{Lift} = 1,16$) và Bìa hồ sơ (\texttt{Binders}, $\text{Confidence} = 23,1\%, \text{Lift} = 1,08$).

\begin{table}[H]
\centering
\caption{Kết quả khai phá luật kết hợp ở cấp độ Tiểu ngành hàng (\texttt{sub\_category}) trên 25.035 giỏ hàng}
\label{tab:assoc_subcat}
\footnotesize
\renewcommand{\arraystretch}{1.22}
\begin{tabularx}{\textwidth}{|c|l|l|r|r|r|X|}
\hline
\textbf{Hạng} & \textbf{Nhóm hàng mua trước ($A$)} & \textbf{Nhóm hàng mua kèm ($B$)} & \textbf{Support} & \textbf{Confidence} & \textbf{Lift} & \textbf{Ý nghĩa nghiệp vụ bán chéo} \\
\hline
1 & \texttt{Copiers} & \texttt{Labels} & 0,98\% & 11,60\% & 1,17 & Gợi ý vật tư nhãn in khi mua máy photocopy. \\
\hline
2 & \texttt{Appliances} & \texttt{Paper} & 1,01\% & 15,00\% & 1,16 & Đề xuất giấy văn phòng khi mua thiết bị. \\
\hline
3 & \texttt{Copiers} & \texttt{Bookcases} & 0,87\% & 10,30\% & 1,13 & Mua sắm đồng bộ thiết bị và nội thất lưu trữ. \\
\hline
4 & \texttt{Machines} & \texttt{Bookcases} & 0,58\% & 10,20\% & 1,12 & Kết hợp máy móc văn phòng và kệ tủ hồ sơ. \\
\hline
5 & \texttt{Copiers} & \texttt{Envelopes} & 0,87\% & 10,30\% & 1,11 & Mua kèm phong bì ấn phẩm cho văn phòng. \\
\hline
6 & \texttt{Machines} & \texttt{Labels} & 0,62\% & 10,80\% & 1,10 & Bán kèm nhãn dán cho thiết bị phần cứng. \\
\hline
7 & \texttt{Appliances} & \texttt{Binders} & 1,56\% & 23,10\% & 1,08 & Gói vật tư lưu trữ văn phòng phổ biến nhất. \\
\hline
8 & \texttt{Accessories} & \texttt{Phones} & 1,55\% & 13,40\% & 1,07 & Bán chéo phụ kiện công nghệ cùng điện thoại. \\
\hline
\end{tabularx}
\end{table}

Ở cấp độ thứ hai -- cấp độ Sản phẩm chi tiết (\texttt{--level product --min-count 3 --min-confidence 0.10 --min-lift 1.2}), hệ thống phát hiện được 27 luật kết hợp chặt chẽ giữa các mặt hàng cụ thể với hệ số Lift rất cao (Bảng \ref{tab:assoc_product}). Đáng chú ý nhất là cặp sản phẩm túi gửi thư \texttt{Jiffy Mailers, Set of 50} và ổ nhớ \texttt{Memorex Flash Drive, USB} đạt độ tin cậy 21,43\% và hệ số $\text{Lift} = 315,57$, hay cặp máy fax \texttt{Brother Fax and Copier, Color} và sổ ghi lời nhắn \texttt{SanDisk Message Books, Recycled} đạt $\text{Lift} = 223,53$. Các luật này phản ánh chính xác thói quen mua sắm theo bộ dụng cụ văn phòng trọn gói của nhóm khách hàng doanh nghiệp.

\begin{table}[H]
\centering
\caption{Top các luật kết hợp mạnh nhất ở cấp độ Sản phẩm SKU chi tiết lưu trong \texttt{dss.association\_rule}}
\label{tab:assoc_product}
\footnotesize
\renewcommand{\arraystretch}{1.22}
\begin{tabularx}{\textwidth}{|c|X|X|r|r|r|}
\hline
\textbf{Hạng} & \textbf{Sản phẩm mua trước (Antecedent $A$)} & \textbf{Sản phẩm gợi ý mua kèm (Consequent $B$)} & \textbf{Support} & \textbf{Confidence} & \textbf{Lift} \\
\hline
1 & \texttt{Jiffy Mailers, Set of 50} & \texttt{Memorex Flash Drive, USB} & 0,012\% & 21,43\% & 315,57 \\
\hline
2 & \texttt{Memorex Flash Drive, USB} & \texttt{Jiffy Mailers, Set of 50} & 0,012\% & 17,65\% & 315,57 \\
\hline
3 & \texttt{SanDisk Message Books, Recycled} & \texttt{Brother Fax and Copier, Color} & 0,012\% & 21,43\% & 223,53 \\
\hline
4 & \texttt{Brother Fax and Copier, Color} & \texttt{SanDisk Message Books, Recycled} & 0,012\% & 12,50\% & 223,53 \\
\hline
5 & \texttt{OIC Staples, Assorted Sizes} & \texttt{Kraft Clasp Envelope, clear poly} & 0,012\% & 12,00\% & 107,29 \\
\hline
6 & \texttt{Hoover Microwave, White} & \texttt{Fellowes Folders, Blue} & 0,012\% & 16,67\% & 106,99 \\
\hline
7 & \texttt{Tenex Light Bulb, Duo Pack} & \texttt{Binney \& Smith Highlighters} & 0,012\% & 16,67\% & 104,31 \\
\hline
8 & \texttt{Apple Office Telephone, Caller ID} & \texttt{Sharp Copy Machine, Color} & 0,012\% & 13,04\% & 96,04 \\
\hline
\end{tabularx}
\end{table}

\subsection{Mô phỏng kịch bản What-If (Price Elasticity Simulation)}

\subsubsection{Bài toán}

Như đã phân tích trong phần khám phá dữ liệu, việc lạm dụng chiết khấu là nguyên nhân hàng đầu gây thất thoát lợi nhuận tại Global Superstore. Trước khi phê duyệt một chương trình điều chỉnh tỷ lệ chiết khấu thêm một lượng $\Delta d$ (có thể tăng chiết khấu $\Delta d > 0$ để kích cầu hoặc cắt giảm chiết khấu $\Delta d < 0$ để bảo vệ biên lợi nhuận) cho một khu vực thị trường và ngành hàng cụ thể, nhà quản trị cần biết trước doanh thu và lợi nhuận ròng sẽ biến đổi như thế nào dưới tác động của quy luật cung cầu.

\subsubsection{Phương pháp}

Mô-đun mô phỏng kịch bản What-If được xây dựng trong \texttt{src/dss/what\_if/simulate.py}. Khi người dùng chọn một thị trường mục tiêu (\texttt{market}) và một ngành hàng (\texttt{category}), hàm \texttt{fetch\_baseline\_from\_db()} lập tức truy vấn trực tiếp vào bảng sự kiện \texttt{gold.fact\_sales} kết nối cùng \texttt{dim\_geography} và \texttt{dim\_product} để trích xuất ba tham số trạng thái gốc thực tế (Baseline) gồm: doanh thu hiện tại $S_{\text{base}}$, lợi nhuận ròng hiện tại $P_{\text{base}}$ và mức chiết khấu trung bình hiện tại $d_{\text{base}}$. Khi áp dụng mức điều chỉnh chiết khấu $\Delta d$, tỷ lệ chiết khấu mới được giới hạn trong miền hợp lệ $d_{\text{new}} = \min(0,85; \max(0,0; d_{\text{base}} + \Delta d))$. Tỷ lệ phần trăm thay đổi của mức giá bán thực tế mà khách hàng phải trả là:
\begin{equation}
    \frac{\Delta P}{P} = \frac{(1 - d_{\text{new}}) - (1 - d_{\text{base}})}{1 - d_{\text{base}}}
\end{equation}
Gọi $e > 0$ là hệ số co giãn của cầu theo giá (Price Elasticity of Demand, mặc định $e = 1,4$ đối với các mặt hàng thương mại điện tử có tính cạnh tranh cao), tỷ lệ thay đổi sản lượng tiêu thụ dự kiến $\frac{\Delta Q}{Q}$, doanh thu mô phỏng $S_{\text{new}}$, tổng chi phí khả biến tương ứng $C_{\text{new}}$ và lợi nhuận ròng mô phỏng $P_{\text{new}}$ được tính toán theo hệ phương trình kinh tế lượng:
\begin{equation}
    \frac{\Delta Q}{Q} = -e \cdot \frac{\Delta P}{P}, \quad
    S_{\text{new}} = S_{\text{base}} \cdot \left(1 + \frac{\Delta Q}{Q}\right) \cdot \frac{1 - d_{\text{new}}}{1 - d_{\text{base}}}
\end{equation}
\begin{equation}
    C_{\text{base}} = S_{\text{base}} - P_{\text{base}}, \quad
    C_{\text{new}} = C_{\text{base}} \cdot \left(1 + \frac{\Delta Q}{Q}\right), \quad
    P_{\text{new}} = S_{\text{new}} - C_{\text{new}}
\end{equation}
Dựa trên kết quả mô phỏng, hệ thống áp dụng bộ quy tắc suy diễn khuyến nghị tự động gồm ba mức: nếu $P_{\text{new}} < 0$ hoặc biên lợi nhuận mới $M_{\text{new}} = (P_{\text{new}}/S_{\text{new}})\times 100\% < 5\%$, hệ thống trả về trạng thái \texttt{[!] WARNING / REJECT} (bác bỏ kịch bản do xói mòn lợi nhuận nghiêm trọng); nếu $M_{\text{new}} \ge 5\%$ nhưng $P_{\text{new}} < P_{\text{base}}$, hệ thống trả về \texttt{[*] CAUTION} (cảnh báo doanh số tăng không bù đắp được chi phí giảm giá); ngược lại nếu $P_{\text{new}} \ge P_{\text{base}}$ và $M_{\text{new}} \ge 5\%$, hệ thống trả về \texttt{[+] APPROVE / VIABLE} (khuyến nghị phê duyệt thực thi).

\subsubsection{Thực nghiệm và kết quả}

Bảng \ref{tab:what_if_results} trình bày kết quả thực nghiệm mô phỏng trên dữ liệu gốc thực tế của ngành hàng Nội thất (\texttt{Furniture}) tại thị trường \texttt{APAC} ($S_{\text{base}} = 1.344.282,66\text{ USD}$, $P_{\text{base}} = 109.126,65\text{ USD}$, $d_{\text{base}} = 15,82\%$, biên lợi nhuận gốc $M_{\text{base}} = 8,12\%$) với hệ số co giãn $e = 1,4$. Ở kịch bản 1, khi tăng chiết khấu thêm $+5,0\%$ (lên $20,82\%$), mặc dù sản lượng cầu tăng giúp doanh thu tăng nhẹ $+2,02\%$ lên 1.371.400,25 USD, nhưng tổng chi phí khả biến tăng theo sản lượng khiến lợi nhuận ròng sụt giảm tới $-68,89\%$ (chỉ còn 33.949,20 USD, biên lợi nhuận rơi xuống 2,48\%), do đó hệ thống tự động bác bỏ (\texttt{REJECT}). Ngược lại, ở kịch bản 2, khi thực hiện cắt giảm $-5,0\%$ chiết khấu (xuống còn $10,82\%$), dù doanh thu giảm nhẹ $-2,99\%$ xuống 1.304.063,11 USD, lợi nhuận ròng lại tăng vọt $+62,67\%$ lên mức 177.515,43 USD và nâng biên lợi nhuận lên 13,61\%, được hệ thống khuyến nghị phê duyệt (\texttt{APPROVE}).

\begin{table}[H]
\centering
\caption{Kết quả thực nghiệm mô phỏng What-If đối với ngành hàng Furniture tại thị trường APAC ($e=1,4$)}
\label{tab:what_if_results}
\footnotesize
\renewcommand{\arraystretch}{1.25}
\begin{tabularx}{\textwidth}{|l|r|r|r|r|X|}
\hline
\textbf{Kịch bản mô phỏng} & \textbf{Chiết khấu} & \textbf{Doanh thu (USD)} & \textbf{Lợi nhuận (USD)} & \textbf{Biên LN} & \textbf{Quyết định tự động của DSS} \\
\hline
Hiện trạng (Baseline CSDL) & 15,82\% & 1.344.282,66 & 109.126,65 & 8,12\% & Trạng thái tham chiếu hiện hành từ tầng Gold. \\
\hline
Kịch bản 1: Tăng CK $+5,0\%$ & 20,82\% & 1.371.400,25 ($+2,02\%$) & 33.949,20 ($-68,89\%$) & 2,48\% & \textbf{\texttt{[!] REJECT:}} Xói mòn lợi nhuận dưới ngưỡng an toàn 5\%. \\
\hline
Kịch bản 2: Giảm CK $-5,0\%$ & 10,82\% & 1.304.063,11 ($-2,99\%$) & 177.515,43 ($+62,67\%$) & 13,61\% & \textbf{\texttt{[+] APPROVE:}} Tối ưu hóa lợi nhuận ròng thêm $+68.388,78$ USD. \\
\hline
\end{tabularx}
\end{table}

\subsection{Tích hợp kết quả vào hệ thống hỗ trợ quyết định}

Toàn bộ kết quả từ các mô-đun phân tích được tích hợp thành một luồng khép kín từ cơ sở dữ liệu lên giao diện người dùng cuối. Cụ thể, sau khi các kịch bản \texttt{src.dss.predictive.run} và \texttt{src.dss.association.run} hoàn tất tính toán, các điểm dự báo chuỗi thời gian và các luật kết hợp giỏ hàng được lưu trữ trực tiếp vào các bảng \texttt{dss.forecast\_result} và \texttt{dss.association\_rule} trong PostgreSQL. Ở tầng dịch vụ, ứng dụng FastAPI (\texttt{src/app/api/main.py}) cung cấp bốn điểm cuối RESTful chuẩn JSON gồm \texttt{GET /api/kpi/summary}, \texttt{POST /api/dss/what-if}, \texttt{GET /api/dss/forecast} và \texttt{GET /api/dss/association-rules} có kiểm tra hợp lệ dữ liệu đầu vào bằng Pydantic.

Ở tầng giao diện điều hành trực tiếp, ứng dụng Streamlit (\texttt{src/app/main.py}) sử dụng cơ chế bộ nhớ đệm \texttt{@st.cache\_data(ttl=60)} để truy vấn trực tiếp dữ liệu từ lược đồ \texttt{gold} và \texttt{dss} lên bốn phân trang chuyên biệt: Phân trang 1 hiển thị bảng điểm KPI toàn cầu, biểu đồ trục kép Doanh thu -- Lợi nhuận qua 48 tháng và cảnh báo vùng rủi ro; Phân trang 2 cung cấp giao diện tương tác thời gian thực với các thanh trượt điều chỉnh chiết khấu và hệ số co giãn cầu $e$, tự động nạp số liệu gốc từ PostgreSQL theo từng thị trường/ngành hàng và hiển thị biểu đồ đối chiếu cùng khuyến nghị phê duyệt/từ chối; Phân trang 3 vẽ đồ thị chuỗi thời gian 48 tháng lịch sử nối tiếp với 6 tháng dự báo tương lai kèm vùng tô màu khoảng tin cậy 95\%; và Phân trang 4 cung cấp công cụ tra cứu, lọc động các luật kết hợp theo ngưỡng Lift để gợi ý tạo gói combo khuyến mãi ngay tại giỏ hàng. Song song đó, tiện ích \texttt{src/app/export\_for\_powerbi.py} và bộ mã nguồn DAX/Power Query M trong \texttt{reports/powerbi/} cho phép đồng bộ toàn bộ 6 bảng tầng Gold và 2 bảng tầng DSS sang báo cáo quản trị Power BI Desktop.

\newpage
\section{Kết quả thực nghiệm và đánh giá}

\subsection{Kết quả xây dựng hệ thống}

\subsubsection{Data Warehouse}

Hệ thống Kho dữ liệu Global Superstore đã được triển khai và vận hành ổn định trên nền tảng PostgreSQL 17 đóng gói trong Docker container, khởi tạo đầy đủ 4 lược đồ chuyên biệt (\texttt{bronze}, \texttt{silver}, \texttt{gold}, \texttt{dss}) với tổng cộng 15 bảng vật lý và 3 bảng tổng hợp Materialized Views thông qua lệnh điều phối duy nhất \texttt{python -m src.pipeline.cli --init-db}. Tại tầng Bronze, hệ thống đã tiếp nhận toàn vẹn 51.290 bản ghi giao dịch gốc vào bảng \texttt{bronze.raw\_orders}, trích xuất chuẩn xác 632 bản ghi mã bưu chính vào \texttt{bronze.postal\_code} và 149 ánh xạ quốc gia -- thị trường vào \texttt{bronze.market}, đồng thời lưu vết kiểm toán SHA-256 vào \texttt{bronze.ingestion}. Tại tầng Silver, bảng \texttt{silver.orders} đã chuẩn hóa thành công 100\% bản ghi sang kiểu dữ liệu định lượng và ngày tháng chuẩn ISO mà không phát sinh bất kỳ lỗi vi phạm ràng buộc miền giá trị nào. Tại tầng Gold, mô hình hình sao hoàn chỉnh đã được thiết lập với bảng sự kiện trung tâm \texttt{gold.fact\_sales} (51.290 dòng) kết nối chặt chẽ qua các khóa thay thế số nguyên với 5 bảng chiều gồm \texttt{gold.dim\_date} (2.922 ngày lịch), \texttt{gold.dim\_customer} (1.590 khách hàng), \texttt{gold.dim\_product} (10.292 sản phẩm), \texttt{gold.dim\_geography} (3.847 vị trí địa lý) và \texttt{gold.dim\_ship\_mode} (4 phương thức giao vận), cùng 3 Data Marts tính toán sẵn (\texttt{mv\_sales\_mart}, \texttt{mv\_product\_mart}, \texttt{mv\_customer\_mart}).

\subsubsection{Dashboard và giao diện truy vấn}

Hệ thống trình diễn và truy vấn dữ liệu được hiện thực hóa song song trên hai nền tảng nhằm phục vụ đa dạng nhu cầu của người sử dụng. Trên nền tảng ứng dụng Web, giao diện điều hành tương tác được xây dựng bằng ngôn ngữ Python và thư viện Streamlit (\texttt{src/app/main.py}) kết hợp cùng dịch vụ RESTful API trên nền tảng FastAPI (\texttt{src/app/api/main.py}). Giao diện Web được chia thành bốn phân trang trực quan kết nối trực tiếp với cơ sở dữ liệu PostgreSQL: Phân trang tổng quan hiển thị 5 thẻ chỉ số KPI toàn cục, hệ thống cảnh báo điểm nóng rủi ro và biểu đồ kết hợp cột -- đường trục kép phản ánh diễn biến Doanh thu, Lợi nhuận và Số lượng đơn hàng qua 48 tháng liên tục; các phân trang còn lại cung cấp công cụ tương tác thời gian thực với ba mô-đun của hệ hỗ trợ quyết định.

Song song với ứng dụng Web, bộ tài nguyên báo cáo quản trị trên phần mềm Power BI Desktop đã được xây dựng hoàn chỉnh trong thư mục \texttt{reports/powerbi/}, bao gồm tệp cấu hình giao diện màu sắc chuẩn nhận diện (\texttt{hcmut\_theme.json}), tập lệnh kết nối trực tiếp cơ sở dữ liệu Power Query M (\texttt{02\_POWER\_QUERY\_M\_SCRIPTS.m}), thư viện hơn 20 công thức đo lường động DAX (\texttt{03\_DAX\_MEASURES.dax}) như \texttt{Total Sales}, \texttt{Total Profit}, \texttt{Profit Margin \%}, \texttt{Sales YoY Growth \%}, \texttt{Pareto \%}, \texttt{Negative Profit Orders Count} và tiện ích xuất dữ liệu sạch (\texttt{src/app/export\_for\_powerbi.py}). Báo cáo được tổ chức thành ba trang phân tích chuyên sâu: Trang 1 (\textit{Executive Sales Performance}) theo dõi hiệu suất tài chính theo thời gian và bản đồ địa lý toàn cầu; Trang 2 (\textit{Product Portfolio \& Pareto ABC}) phân tích ma trận định vị sản phẩm và quy luật Pareto 80/20; Trang 3 (\textit{Customer Analytics \& Cross-Selling}) phân tích cơ cấu phân khúc khách hàng và bảng đề xuất bán chéo sản phẩm.

\subsubsection{Hệ thống hỗ trợ quyết định (DSS)}

Phân hệ DSS đã được hiện thực hóa đầy đủ với ba mô-đun thuật toán hoạt động độc lập qua giao diện dòng lệnh (CLI) hoặc gọi trực tiếp từ giao diện Web và API. Mô-đun dự báo chuỗi thời gian (\texttt{src/dss/predictive/run.py}) đã huấn luyện thành công mô hình San bằng mũ có xu thế trên chuỗi 48 tháng lịch sử, lưu trữ kết quả đánh giá vào bảng \texttt{dss.forecast\_run} và xuất 6 điểm dự báo doanh thu nửa đầu năm 2015 kèm cận dưới/cận trên của khoảng tin cậy 95\% vào bảng \texttt{dss.forecast\_result}. Mô-đun khai phá giỏ hàng (\texttt{src/dss/association/run.py}) đã phân tích 25.035 đơn đặt hàng và lưu trữ bền vững vào bảng \texttt{dss.association\_rule} cả hai cấp độ tri thức: 14 luật kết hợp ở cấp độ tiểu ngành hàng và 27 luật kết hợp ở cấp độ sản phẩm SKU chi tiết. Mô-đun mô phỏng kịch bản What-If (\texttt{src/dss/what\_if/simulate.py}) đã tích hợp hoàn chỉnh khả năng truy vấn động số liệu gốc từ kho dữ liệu Gold theo từng bộ lọc thị trường -- ngành hàng, giải hệ phương trình co giãn cầu theo giá trong thời gian thực và đưa ra thông điệp khuyến nghị hành động ba mức độ rõ ràng.

\subsection{Đánh giá Data Warehouse và ETL}

Để đánh giá toàn diện chất lượng của Kho dữ liệu và quy trình ETL, nhóm thực hiện xem xét trên bốn tiêu chí kỹ thuật cốt lõi gồm tính chính xác dữ liệu, tính lũy đẳng và khả năng nạp gia tăng, hiệu năng xử lý đường ống và tốc độ phản hồi truy vấn OLAP.

Về tính chính xác và toàn vẹn dữ liệu, kết quả thực thi bộ kiểm định tự động tại cả hai tầng Silver và Gold đạt tỷ lệ vượt qua 100\% trên toàn bộ các bài kiểm tra. Không có bất kỳ bản ghi nào bị rỗng khóa bắt buộc, không có bản ghi trùng lặp \texttt{row\_id} và không tồn tại khóa ngoại mồ côi giữa bảng sự kiện \texttt{gold.fact\_sales} với năm bảng chiều. Phép đối soát tổng hợp tài chính giữa tầng Silver và tầng Gold khớp nhau tuyệt đối trên cả 51.290 bản ghi với tổng doanh thu đạt $12.642.501,91\text{ USD}$, tổng lợi nhuận đạt $1.467.457,29\text{ USD}$ và tổng số lượng hàng hóa đạt $178.312\text{ đơn vị}$. Các ca biên phức tạp của tập dữ liệu gốc như ngày tháng đa định dạng, 80,51\% mã bưu chính khuyết thiếu và hiện tượng thành phố thuộc nhiều khu vực thị trường đều đã được xử lý triệt để mà không phải loại bỏ bất kỳ giao dịch hợp lệ nào.

Về tính lũy đẳng (Idempotency) và khả năng nạp gia tăng (Incremental Loading), hệ thống chứng tỏ độ tin cậy cao trong vận hành thực tế. Nhờ cơ chế đối chiếu mã băm SHA-256 trong bảng \texttt{bronze.ingestion} kết hợp mệnh đề \texttt{ON CONFLICT DO NOTHING / DO UPDATE} tại tầng Silver và Gold, việc kích hoạt lại đường ống ETL nhiều lần trên cùng một tệp dữ liệu hoàn toàn không làm phát sinh bản ghi dư thừa. Khi kiểm thử nạp gia tăng với lô 500 bản ghi mới của Quý 1/2015, cơ chế cờ đánh dấu \texttt{transform = FALSE} giúp hệ thống chỉ xử lý đúng 500 bản ghi mới phát sinh trong thời gian chưa đầy 3 giây, nâng tổng số dòng trong \texttt{gold.fact\_sales} lên 51.790 dòng một cách chính xác.

Về hiệu năng truy vấn phân tích OLAP, việc chuyển đổi từ bảng phẳng sang mô hình hình sao sử dụng khóa thay thế số nguyên kết hợp đánh chỉ mục B-Tree trên toàn bộ các cột khóa ngoại đã mang lại tốc độ truy vấn vượt trội. Toàn bộ 7 truy vấn OLAP phức tạp (bao gồm các phép toán nặng như \texttt{GROUP BY ROLLUP} trên 4 chiều, \texttt{GROUP BY CUBE} sinh 8 tổ hợp chiều và các hàm cửa sổ \texttt{DENSE\_RANK}, \texttt{LAG}) khi chạy trực tiếp trên 51.290 dòng của bảng sự kiện đều hoàn tất trong khoảng thời gian từ $16,1\text{ ms}$ đến $49,2\text{ ms}$ (trung bình dưới $30\text{ ms}$). Đặc biệt, khi truy vấn thông qua 3 bảng Materialized Views đã được tổng hợp trước, thời gian phản hồi giảm xuống dưới $5\text{ ms}$, đảm bảo trải nghiệm mượt mà tức thì khi người dùng thao tác lọc chéo trên giao diện Web và Power BI.

\subsection{Đánh giá DSS}

Tính hiệu quả của Hệ hỗ trợ quyết định được đánh giá dựa trên độ chính xác của các mô hình định lượng và giá trị thực tiễn của các tri thức được khám phá đối với hoạt động quản trị tại Global Superstore. Đối với mô-đun dự báo chuỗi thời gian (Predictive Analytics), việc phân tách tập huấn luyện (2011--2013) và tập kiểm thử (2014) theo đúng trình tự thời gian giúp đánh giá trung thực năng lực dự báo. Mặc dù doanh số năm 2014 có biên độ dao động mùa vụ rất lớn giữa các tháng đầu năm và cuối năm, mô hình San bằng mũ kết hợp xu thế tuyến tính vẫn bám sát quỹ đạo tăng trưởng chung và đạt sai số phần trăm tuyệt đối trung bình $\text{MAPE} = 14,72\%$ ($\text{MAE} = 48.591,40\text{ USD}$). So với mô hình trung bình trượt cơ sở vốn bị trễ pha khi chuỗi thời gian có xu thế dốc lên, mô hình được chọn phản ánh chính xác tốc độ tăng trưởng trung bình hàng tháng và cung cấp dải tin cậy 95\% hợp lý cho 6 tháng đầu năm 2015.

Đối với mô-đun khai phá luật kết hợp (Association Rule Mining), việc thiết kế chiến lược khai phá hai cấp độ đã giải quyết triệt để bài toán thưa dữ liệu của danh mục 10.292 sản phẩm. Ở cấp độ tiểu ngành hàng (\texttt{sub\_category}), 14 luật tìm được đều có độ hỗ trợ từ 0,58\% đến 1,56\% và độ tin cậy lên tới 23,10\%, cung cấp cơ sở vững chắc để bố trí các gian hàng liên quan trên trang thương mại điện tử (như gợi ý Nhãn dán, Phong bì và Tủ sách ngay khi khách hàng xem danh mục Máy sao chụp). Ở cấp độ sản phẩm chi tiết (\texttt{product}), việc sử dụng ngưỡng đếm tuyệt đối $\text{min\_count} = 3$ đã giúp phát hiện 27 cặp sản phẩm có hệ số Lift rất cao (từ $31,91$ đến $315,57$), phục vụ trực tiếp cho tính năng đề xuất mua kèm tự động tại bước kiểm tra giỏ hàng.

Đối với mô-đun mô phỏng kịch bản What-If, giá trị lớn nhất mà mô-đun mang lại là việc kết nối trực tiếp lý thuyết độ co giãn cầu theo giá với số liệu biên lợi nhuận thực tế của từng thị trường trong kho dữ liệu. Thực nghiệm trên ngành hàng Nội thất và tiểu ngành hàng Bàn làm việc (\texttt{Tables}) đã chứng minh định lượng vì sao việc tăng chiết khấu tại các nhóm hàng có biên lợi nhuận gốc mỏng ($< 10\%$) lại gây ra sự sụt giảm lợi nhuận ròng nghiêm trọng dù doanh số bán hàng vẫn tăng, từ đó giúp ban giám đốc tránh được các chiến dịch khuyến mãi gây lỗ vốn.

\subsection{Hạn chế và hướng phát triển}

Bên cạnh những kết quả đã đạt được, hệ thống hiện tại vẫn tồn tại một số hạn chế nhất định xuất phát từ đặc thù của tập dữ liệu nguồn và phạm vi triển khai trong khuôn khổ môn học. Thứ nhất, dữ liệu lịch sử của Global Superstore chỉ bao gồm các giao dịch bán hàng đã hoàn tất mà không chứa dữ liệu về số lượng hàng tồn kho thực tế tại từng kho bãi (Inventory on hand), giá vốn nhập hàng chi tiết theo từng nhà cung cấp hoặc chi phí tiếp thị quảng cáo (Marketing Spend), do đó hệ thống chưa thể mở rộng sang các bài toán tối ưu hóa điểm đặt hàng lại (Reorder Point -- ROP) hay tính toán chi phí thu hút khách hàng (CAC). Thứ hai, ở mô-đun dự báo chuỗi thời gian, công thức dự báo điểm ở các bước xa ($h > 1$) hiện đang ngoại suy theo độ dốc xu thế tuyến tính kết hợp làm mượt mũ đơn giản nên chưa mô hình hóa đầy đủ thành phần dao động hình sin của chu kỳ mùa vụ 12 tháng vào các điểm dự báo xa. Thứ ba, quy trình nạp dữ liệu hiện tại vẫn vận hành theo phương thức xử lý theo lô (Batch Processing) kích hoạt qua dòng lệnh thay vì tự động lắng nghe sự kiện theo luồng thời gian thực.

Từ những hạn chế trên, nhóm đề xuất bốn hướng phát triển cụ thể cho hệ thống trong tương lai. Hướng phát triển thứ nhất là tích hợp công cụ điều phối quy trình công việc tự động như Apache Airflow hoặc Prefect để lập lịch chạy pipeline nạp gia tăng định kỳ hàng đêm kèm theo cơ chế gửi cảnh báo tự động khi phát hiện dữ liệu bất thường. Hướng phát triển thứ hai là nâng cấp mô-đun dự báo chuỗi thời gian bằng cách tích hợp các mô hình có thành phần mùa vụ tường minh như SARIMA (Seasonal ARIMA), Holt-Winters ba thành phần (Triple Exponential Smoothing) hoặc mô hình học máy Gradient Boosting (XGBoost, LightGBM) với các đặc trưng trễ thời gian (Lag features) để giảm sai số MAPE xuống dưới 8\%. Hướng phát triển thứ ba là mở rộng chiến lược Chiều biến đổi chậm tại bảng \texttt{gold.dim\_customer} và \texttt{gold.dim\_product} từ SCD Type 1 sang SCD Type 2 hoàn chỉnh với các cột \texttt{effective\_from}, \texttt{effective\_to} và \texttt{is\_current} nhằm theo dõi lịch sử chuyển đổi phân khúc của khách hàng qua nhiều năm. Hướng phát triển thứ tư là kết hợp kết quả phân khúc khách hàng RFM tại \texttt{gold.mv\_customer\_mart} với các luật kết hợp giỏ hàng để xây dựng hệ thống gợi ý cá nhân hóa (Personalized Recommendation System) cho từng nhóm khách hàng mục tiêu.

\newpage
\section{Kết luận}

\subsection{Tổng kết kết quả đạt được}

Báo cáo Bài tập lớn môn Kho dữ liệu và Hệ hỗ trợ quyết định (CO4031) đã trình bày toàn bộ quá trình nghiên cứu, phân tích nghiệp vụ, thiết kế kiến trúc và hiện thực hóa hệ thống Kho dữ liệu tích hợp Hệ hỗ trợ quyết định cho bài toán doanh nghiệp bán lẻ toàn cầu Global Superstore. Đối chiếu với các mục tiêu đề ra ban đầu, nhóm thực hiện đã hoàn thành trọn vẹn toàn bộ các hạng mục kỹ thuật và nghiệp vụ từ tầng dữ liệu thô đến tầng giao diện ra quyết định.

Thứ nhất, nhóm đã tiến hành khảo sát sâu bộ dữ liệu 51.290 bản ghi giao dịch trên 147 quốc gia trong giai đoạn 2011--2014, xác lập hệ thống 13 chỉ số đo lường hiệu quả kinh doanh (KPIs) chuẩn hóa và chỉ ra các vấn đề cốt lõi về chất lượng dữ liệu nguồn cũng như bài toán thất thoát lợi nhuận do chính sách chiết khấu tại khu vực EMEA và tiểu ngành hàng Bàn làm việc (\texttt{Tables}).

Thứ hai, nhóm đã thiết kế và triển khai thành công Kho dữ liệu theo kiến trúc Medallion ba tầng (\texttt{bronze}, \texttt{silver}, \texttt{gold}) trên hệ quản trị cơ sở dữ liệu PostgreSQL 17 kết hợp Docker. Quy trình ETL được tự động hóa hoàn toàn, tích hợp cơ chế kiểm toán mã băm SHA-256, xử lý triệt để đa định dạng ngày tháng và dữ liệu mã bưu chính khuyết thiếu, đảm bảo tính lũy đẳng và hỗ trợ nạp gia tăng (Incremental Loading) thông qua cơ chế bắt thay đổi dữ liệu bằng cờ trạng thái \texttt{transform}. Mô hình hình sao (Star Schema) tại tầng Gold với 1 bảng sự kiện (\texttt{gold.fact\_sales}) và 5 bảng chiều sử dụng khóa thay thế số nguyên đã vượt qua 100\% kịch bản kiểm định chất lượng và đối soát số liệu tài chính tuyệt đối ($12.642.501,91\text{ USD}$ doanh thu, $1.467.457,29\text{ USD}$ lợi nhuận). Hệ thống 7 truy vấn phân tích OLAP và 3 bảng Materialized Views (\texttt{mv\_sales\_mart}, \texttt{mv\_product\_mart}, \texttt{mv\_customer\_mart}) đạt thời gian phản hồi thực tế dưới $50\text{ ms}$.

Thứ ba, nhóm đã phát triển hoàn chỉnh Hệ hỗ trợ quyết định (DSS) gồm ba mô-đun định lượng kết nối trực tiếp với kho dữ liệu: mô-đun dự báo chuỗi thời gian bằng thuật toán San bằng mũ có xu thế dự báo doanh thu 6 tháng đầu năm 2015 kèm khoảng tin cậy 95\% (đạt $\text{MAPE} = 14,72\%$ trên tập kiểm thử); mô-đun khai phá luật kết hợp Apriori hai cấp độ trên 25.035 giỏ hàng phát hiện 14 luật cấp tiểu ngành hàng và 27 luật cấp sản phẩm SKU (với hệ số Lift lên tới $315,57$); và mô-đun mô phỏng kịch bản What-If dựa trên độ co giãn của cầu theo giá. Toàn bộ kết quả phân tích được tích hợp đồng bộ lên dịch vụ RESTful API (FastAPI), ứng dụng Web tương tác 4 phân trang (Streamlit) và bộ báo cáo quản trị 3 trang trên Power BI Desktop, tạo thành một giải pháp hỗ trợ ra quyết định hoàn chỉnh, trực quan và có giá trị ứng dụng thực tiễn cao.

\subsection{Định hướng phát triển}

Trong giai đoạn tiếp theo, hệ thống có tiềm năng mở rộng mạnh mẽ theo hướng kiến trúc dữ liệu hiện đại quy mô doanh nghiệp lớn. Về mặt hạ tầng và kỹ thuật dữ liệu, việc tích hợp nền tảng điều phối luồng công việc tự động (Apache Airflow) và cơ chế bắt thay đổi dữ liệu dựa trên nhật ký giao dịch (Log-based CDC với Debezium và Apache Kafka) sẽ giúp chuyển đổi quy trình nạp dữ liệu từ xử lý theo lô định kỳ sang xử lý gần thời gian thực (Near Real-time Data Warehousing). Bên cạnh đó, việc nâng cấp các bảng chiều khách hàng và sản phẩm lên chuẩn SCD Type 2 hoàn chỉnh và bổ sung thêm các bảng sự kiện quản lý tồn kho (Inventory Snapshot Fact Table) sẽ giúp mở rộng phạm vi phân tích sang toàn bộ chuỗi cung ứng.

Về mặt mô hình hỗ trợ quyết định và trí tuệ nhân tạo, hệ thống có thể được nâng cấp bằng cách huấn luyện các mô hình dự báo chuỗi thời gian phân cấp (Hierarchical Time Series Forecasting) kết hợp các thuật toán học máy như Prophet, SARIMA và LightGBM để dự báo nhu cầu riêng biệt cho từng thị trường và từng nhóm ngành hàng. Đồng thời, sự kết hợp giữa ma trận phân khúc khách hàng RFM đã có sẵn trong \texttt{gold.mv\_customer\_mart} với thuật toán lọc cộng tác (Collaborative Filtering) và các luật kết hợp Apriori sẽ tạo tiền đề xây dựng một bộ máy khuyến nghị sản phẩm cá nhân hóa toàn diện, góp phần tối ưu hóa đồng thời doanh thu bán chéo và biên lợi nhuận ròng cho doanh nghiệp.

\newpage
% ===== Tài liệu tham khảo =====

\begin{thebibliography}{9}
\addcontentsline{toc}{section}{Tài liệu tham khảo}

\bibitem{kimball2013}
Ralph Kimball, Margy Ross (2013), \textit{The Data Warehouse Toolkit: The Definitive Guide to Dimensional Modeling}, 3rd Edition, John Wiley \& Sons, Inc.

\bibitem{turban2014}
Efraim Turban, Ramesh Sharda, Dursun Delen (2014), \textit{Business Intelligence and Analytics: Systems for Decision Support}, 10th Edition, Pearson Education.

\bibitem{simon1960}
Herbert A. Simon (1960), \textit{The New Science of Management Decision}, Harper \& Row, New York.

\bibitem{agrawal1994}
Rakesh Agrawal, Ramakrishnan Srikant (1994), ``Fast Algorithms for Mining Association Rules in Large Databases'', \textit{Proceedings of the 20th International Conference on Very Large Data Bases (VLDB)}, pp. 487--499.

\bibitem{hyndman2021}
Rob J. Hyndman, George Athanasopoulos (2021), \textit{Forecasting: Principles and Practice}, 3rd Edition, OTexts, Melbourne, Australia.

\bibitem{postgresql17}
The PostgreSQL Global Development Group (2024), \textit{PostgreSQL 17 Official Documentation -- Data Warehousing, Window Functions, and Materialized Views}, \url{https://www.postgresql.org/docs/17/}.

\end{thebibliography}

\end{document}